
\documentclass[11pt]{article}
\usepackage[a4paper,margin=1in]{geometry}
\usepackage[T1]{fontenc}
\usepackage[utf8]{inputenc}
\usepackage{lmodern}
\usepackage{amsmath,amssymb,amsthm,mathtools,mathrsfs}
\usepackage{enumitem}
\usepackage{microtype}
\usepackage[hidelinks]{hyperref}
\usepackage{graphicx}
\setlength{\parindent}{1.2em}
\setlength{\parskip}{0.25em}
\allowdisplaybreaks
\sloppy
\DeclareMathOperator{\Spec}{Spec}
\DeclareMathOperator{\Hom}{Hom}
\DeclareMathOperator{\Ext}{Ext}
\DeclareMathOperator{\Pic}{Pic}
\DeclareMathOperator{\Aut}{Aut}
\DeclareMathOperator{\Ob}{Ob}
\DeclareMathOperator{\Supp}{Supp}
\DeclareMathOperator{\rank}{rank}
\newcommand{\cO}{\mathcal O}
\newcommand{\cC}{\mathcal C}
\newcommand{\cX}{\mathcal X}
\newcommand{\cY}{\mathcal Y}
\newcommand{\cZ}{\mathcal Z}
\newcommand{\cM}{\mathcal M}
\newcommand{\Z}{\mathbb Z}
\newcommand{\Q}{\mathbb Q}
\newcommand{\C}{\mathbb C}
\newcommand{\A}{\mathbb A}
\newcommand{\Pj}{\mathbb P}
\theoremstyle{plain}
\newtheorem{theorem}{Theorem}
\newtheorem{proposition}{Proposition}
\newtheorem{lemma}{Lemma}
\newtheorem{corollary}{Corollary}
\theoremstyle{definition}
\newtheorem{definition}{Definition}
\theoremstyle{remark}
\newtheorem{remark}{Remark}
\usepackage[french,english]{babel}

\graphicspath{{assets/}}
\newcommand{\displaycrop}[2][0.82\linewidth]{\par\smallskip\begin{center}\includegraphics[width=#1]{#2}\end{center}\smallskip}

\begin{document}
\begin{center}
{\Large\bfseries L'irr\'eductibilit\'e de l'espace des courbes de genre donn\'e}\\[0.7em]
{\large P. Deligne et D. Mumford}
\end{center}
Fixons un corps alg\'ebriquement clos k. Soit M\_g be the espace de modules of curves of genus g over k. Le r\'esultat principal de cette note est que M\_g est irr\'eductible pour tout k. Bien entendu, le fait que M\_g is irr\'eductible ne d\'epend que de la caract\'eristique of k. When la caract\'eristique is o, nous pouvons supposer that k=C, and then le r\'esultat est classique. Une preuve simple se trouve dans Enriques-Chisini [E, vol. 3, chap. 3], fond\'ee sur l'analyse de the totality of coverings of P\textasciicircum{}1 of degree n, with a fixed number d of ordinary branch points. Cette m\'ethode a \'et\'e \'etendue to char. p by William Fulton [F], en utilisant des sp\'ecialisations from char. 0 to char. p provided that p>2g+1. Malheureusement, les tentatives pour \'etendre cette m\'ethode to all p semblent buter sur des questions difficiles of de ramification sauvage. Aujourd'hui, the Teichm\"uller theory donne une compr\'ehension enti\`erement analytique mais tr\`es profonde into this irreducibility when k=C. Notre approche est cependant la plus proche de Severi's preuve incompl\`ete ([Se], Anhang F; the error is on pp. 344-345 and seems to be quite basic) and suit une suggestion of Grothendieck for using the result in char. 0 pour d\'eduire the result in char. p. La base de both Severi's and Grothendieck's ideas est de construire des familles de courbes X, some singular, with pa(X)-=g, over non-singular parameter spaces, which in some sense contain enough courbes singuli\`eres to link together any two composantes that M\_g might have. Le point essentiel that fait fonctionner cette m\'ethode now is a recent " th\'eor\`eme de r\'eduction stable " for vari\'et\'es ab\'eliennes. This result fut d'abord d\'emontr\'e ind\'ependamment in char. 0 by Grothendieck, au moyen de m\'ethodes de cohomologie \'etale (private correspondence with J. Tate), and by Mumford, applying the easy half of Th\'eor\`eme (2.5), to go from curves to vari\'et\'es ab\'eliennes (cf. [M2]). Grothendieck a r\'ecemment renforc\'e his method de sorte qu'il s'applique en toute caract\'eristique (SGA 7, 1968) 9 Mumford has also given a proof using theta functions in char. \textasciitilde{}2. Le r\'esultat est le suivant:

\begin{quote}Th\'eor\`eme de r\'eduction stable. --- Soit R soit un anneau de valuation discr\`ete with corps des fractions K.\end{quote}
Soit A soit unn abelian variety over K. Then there exists a extension alg\'ebrique finie L of K such (x) The first author wishes to thank the Institut des Hautes ]\textasciitilde{}tudes scientifiques, Bures-sur-Yvette, for support in this resea$\pi$h and N. KATZ, for his invaluable assistance in the preparation of this manuscript; the second author wishes to thank the Tata Institute of Fundamental Resea$\pi$h, Bombay, and the Institut des Hautes ]\textasciitilde{}tudes scientifiques.

that, if R L= cl\^oture int\'egrale of R in L, and if di\textasciitilde{} is the mod\`ele de N\'eron of A $\times$ \textasciitilde{} L over Rs, then the closedfibre AL. \textasciitilde{} of di\textasciitilde{} has no radical unipotent. Nous donnerons deux preuves apparent\'ees of our r\'esultat principal. One of these is quite elementary, and follows by quite techniques standard once the Th\'eor\`eme de r\'eduction stable for vari\'et\'es ab\'eliennes is applied, in w 2, pour d\'eduire an analogous th\'eor\`eme de r\'eduction stable pour les courbes. L'autre preuve is plus puissante, and is based on the use of a cat\'egorie plus grande than the cat\'egorie des sch\'emas, and on proving for the objets de cette cat\'egorie many of the th\'eor\`emes standard for sch\'emas, especially the Enriques-Zariski th\'eor\`eme de connexit\'e (EGA 3, (4-3)). Malheureusement, this cat\'egorie plus grande is pas tout \`a fait une cat\'egorie --- it is a simple type of 2-cat\'egorie; in fact, if X, Y are objets, then H\textasciicircum{}0m(X, Y) is itself a cat\'egorie, but one in which all morphismes are isomorphismees. The objets of this 2-cat\'egorie we call champs alg\'ebriques (1). The espace de modules M\_g is just the " vari\'et\'e grossi\`ere sous-jacente " of a more fundamental objet, the champ des modules =gggstudied in [Ma]. Les d\'etails complets on the basic propreties and theorems for champs alg\'ebriques seront donn\'es ailleurs. In this paper, nous ne donnerons que d\'efinitions and state sans preuve the general theorems which we apply. En utilisant la m\'ethode of champs alg\'ebriques, nous pouvons d\'emontrer non seulement the irreducibility of M\_g itself, mais aussi de tous higher level espace de moduless of curves too (cf. § 5 below).

\section*{1. Courbes stables et leurs modules}
La d\'efinition cl\'e de tout l'article is this:

\begin{quote}D\'efinition (x.x). --- Soit S soit unny sch\'ema. Soit g>2. A stable curve of genus g\end{quote}
over S is a morphismee propre et plat $\pi$:C$\to$S whose fibres g\'eom\'etriques are r\'eduites, connexes, 1-dimensional sch\'emas C8 such that: (i) C8 has only ordinary doublepoints; (ii) /f E is a composante rationnelle non singuli\`ere of Cs, then E meets the other composantes of C\textasciitilde{} in more than 2 points; (iii) dim H\textasciicircum{}1(O\_\{C\_s\})=g. Nous \'etudierons dans cette section trois aspects of la th\'eorie des stable curves: their pluri-canonical syst\`emes lin\'eaires, their d\'eformations, and their automorphismees. Supposons que $\pi$ : C$\to$S is a stable curve. Puisque $\pi$ is plat and its fibres g\'eom\'etriques are intersections compl\`etes locales, the morphisme $\pi$is localement a complete intersection (i.e., localement, C is isomorphe as S-sch\'ema to V(fl, ...,fn\_l)cAn$\times$ where UcS is open, and fl,.-.,f,-l\textasciitilde{}F(0A, xU) are a regular sequence). Par cons\'equent, by la th\'eorie des duality of faisceaux coh\'erents [H], there is a faisceau inversible canonique $\omega$\_c/s on C --- the unique groupe de cohomologie non nul of the complex of sheaves f\textasciicircum{}!(0s). Nous aurons besoin de savoir the faits suivants about \textasciitilde{}c/s: a) pour tous les morphismees f: T$\to$S, \textasciitilde{}CxsW/W is canonically isomorphe to f*(Oc/s) ; (1) A slightlylessgeneral cat\'egorieof objets,called algebraicspaces,has been introduced and studied very deeply by M. ARTIN[Ai] and D. KmrrsoN[K]. The idea of enlarging the cat\'egorie ofvarietiesfor the study of espace de modulessis due originally,we believe,to A. Weil.

b) if S = Spec(k), k alg\'ebriquement clos, let f: C'\textasciitilde{}C be the normalisation of C, xl, ..., xn, yl, ...,y\textasciitilde{} the points of C' such that the zi=f(x\textasciitilde{})=f(yi), i<i<n, are the points doubles of C. Then $\omega$\_c/s is the sheaf of i-forms \textasciitilde{} on C' r\'eguli\`ere sauf pour des p\^oles simples at the x's andy's and with Res\textasciitilde{}i(\textasciitilde{})+Re\%i(\textasciitilde{})=o ; c) if S = Spec(k), and o\textasciitilde{} is a coherent sheaf on C, then H\textasciicircum{}0m(H'(C, o\textasciitilde{}), k) \textasciitilde{} H\textasciicircum{}0m\textasciitilde{}o(o\textasciitilde{}, $\omega$\_e/s).

\begin{quote}Th\'eor\`eme (x .2). --- If g>2 and C is a stable curve of genus g over an alg\'ebriquement clos\end{quote}
corps k, then H\textasciicircum{}1(C, | --- $\omega$\_c/k)---(o) if n>2, and (o| is very ample if n>3. D\'emonstration. --- Puisque C is stable, of genus g>2, every irr\'eductible component E of C either i) has (arithmetic) genus >2 itself, 2) has genus I, but meets other composantes of C in at least one point, or 3) is non-singular, rational and meets other composantes of C in at least three points. But by b) above, toC/k| is isomorphe to toE/k(\textasciitilde{}Q\textasciitilde{}), where \{Q\textasciitilde{}\} are the points where E meets the rest of C. Puisque the degree of toE/kis 2gE---2, it follows that in any of the cases I, 2 or 3, $\omega$\_c/kQg)E has positive degree. This shows immediately that $\omega$\_c/kis ample on each component E of C, hence is ample. l-][0/'aOl---n). Puisque 0E has positive Next, by c) above, Hl($\omega$\_c\textasciitilde{}) is dual to \textasciitilde{} \textasciitilde{}c/k $\omega$\_c/k| l..I0/,.,| 1 --- n'\textasciitilde{} \_\_ degree, \textasciitilde{}C/k'"e$\simeq$"C\textasciitilde{}0E\textasciitilde{}has no sections for any E, any n>2;\_ therefore .\textasciitilde{} \textasciitilde{}"c/k j---(o) if n>2, and so ni($\omega$\_c\textasciitilde{})=(o) if n>2. To prove that an invertible sheaf \textasciitilde{} on any sch\'ema C, propre over k, is very ample, it suffices to show a) for all point ferm\'es x+y H\textasciicircum{}0 5e)$\to$H\textasciicircum{}0 (\textasciitilde{}'|174174 is surjectif, b) for all point ferm\'es x, H\textasciicircum{}0 5r ---$\to$H\textasciicircum{}0 \textasciitilde{}r174 0=/m\textasciitilde{}) is surjectif. Using the exact sequence of cohomology, these both follow if H\textasciicircum{}1(C, m\textasciitilde{}. mu.\textasciitilde{} 9\textasciitilde{} = (o) for all point ferm\'es x, yeC. In our case, s $\omega$\_c| n>3 ' so if we use duality, we must show : (*) H\textasciicircum{}0m(m\textasciitilde{}. mu, to$\simeq$") = (o), if n> 2. If x is a non-singular point, mx is an invertible sheaf. If x is a double point, let $\pi$ : C'$\to$ C be the result of blowing up x, and let x\textasciitilde{}, x2eC' be the two points in $\pi$-a(x). Then it is easy to check that for any invertible sheaf \textasciitilde{}qo on C: H\textasciicircum{}0m(m\textasciitilde{}, s162\textasciitilde{} H\textasciicircum{}0 ', \textasciitilde{}*s H\textasciicircum{}0rn(m2+, s H\textasciicircum{}0(C ', r:*Lf(xt +x2) ). Par cons\'equent, we have 3 cases of (*) to check: Case 1. --- x, y non-singular points of C, then H\textasciicircum{}0 if n>2.

Case 2.- x double point of C, $\pi$:C'$\to$C blowing up x, \{xl, x2\}=$\simeq$a(x), and y a non-singular point of C. Then H\textasciicircum{}0 = (o), H\textasciicircum{}0 if n>2. Case 3. --- x,y points doubles of C, $\pi$ : C'$\to$C blowing up x andy. H\textasciicircum{}0 = (o) if n>2. Then Now since the degree of Ocn (n>2) on all composantes E of C is less than or equal to ---2, all of this is clear, except in those cases where n = 2, the degree of $\omega$\_c on some E is i, and in which two poles are allowed on E. This occurs if: (i) case I, pa(E)=I, E meets C---E at only one point, x, yeE. (ii) case I, pa(E)=o, E meets C---E at only three points, x, yeE. (iii) case 2, E a rational curve with one double point meeting C---E at one point, x=double point of E. But in all these cases, C has composantes besides E and a section in the H \textasciitilde{} in question must definily vanish on all these other composantes. So at the points where E meets C---E, the section has extra zeroes. Puisque the sheaf in question has degree o on E, the section is zero on E too. Q.E.D. Corollaire. --- Soit \textasciitilde{} : C$\to$S soit unny stable curve of genus g>\textasciitilde{}.\_ Then \textasciitilde{}c/s| is relatively very ample if n>3 and \textasciitilde{},(o)\textasciitilde{}) is a localementfree sheqf on S of rank (2n---I)(g---i). D\'emonstration. --- In fact, since for all seS, Hl(o\textasciitilde{}c\textasciitilde{}| it follows from [EGA, chap. 3, w 7], that rL(O)c\textasciitilde{}) is localement free and that \%(\textasciitilde{})|176174 Par cons\'equent the corollary follows. Q.E.D. Taking n = 3, it follows that every stable curve C/S can be realized as a family of curves in psg-6 with H\textasciicircum{}1lbert polynomial: Pa(n) = (6n--- z)(g--- z). ([MI], P. 99), it is easy to prove that there is a Following standard arguments subsch\'ema H\_g r I-Iilb\textasciitilde{}g\_ s of" all " tri-canonically embedded stable curves. (H\textasciicircum{}1lb is the H\textasciicircum{}1lbert sch\'ema over Z.) To be precise, there is an isomorphismee of foncteurs: i set of stable curves \textasciitilde{} : C-.\textasciitilde{}S, plus isomorphismees: 1 H\textasciicircum{}0m(S, H\_g) \textasciitilde{} i P(rL(\textasciitilde{}c\textasciitilde{})) \textasciitilde{} p5g-G x S t (modulo isomorphismee) We will denote by Zg c H\_g $\times$ \textasciitilde{}g-6 the universal tri-canonically embedded stable curve. Le foncteur of stable curves itself is the sheafification of the quotient of foncteurs: H\_g/PGL(5g---6).

We now consider the deformation theory of stable curves. Soit k soit unny ground corps. The deformation theory of X's lisse over k can be found in [SGA, 6o-6I]; for singular S's, the theory has been worked out in [Sc]. We shall indicate here the results of this theory for a sch\'ema X which is (i) one-dimensional; (ii) generically lisse over k; (iii) localement a complete intersection. The advantages of this case are two-fold: first, the " cotangent complex " of Grothendieck, Lichtenbaum and Schlessinger reduces, in view of (ii) and (iii), to the single coherent sheaf fZx/k, the Kahler differentials. Secondly, we have:

\begin{quote}Lemme (I. 3). --- Ext2(f$\Omega$\_X/k, (gx)= (o).\end{quote}
\begin{quote}D\'emonstration. --- Use the spectral sequence: HP(X, Extq(f$\Omega$\_X/k, O\_X)) \textasciitilde{} ExtP+q(ax/k,\textasciitilde{})x)Then (i) He(X, Ext \textasciitilde{} = (o) since dim X = I. (ii) Puisque $\Omega$\_X/k is localement free except at a fini number of points, Extl(flx/k, Ox) has o-dimensional support, hence H\textasciicircum{}1(X, Extl)=(o). (iii) Locally, if we embed XcA n, then f2x/k has a free resolution of length 2: o \textasciitilde{} J/fi $\to$ \textasciitilde{}An| Ox \textasciitilde{} g$\Omega$\_X/k \textasciitilde{} o where 3"= sheaf of ideals defining X. Par cons\'equent Ext2= (o). Q.E.D. In Schlessinger's theory, the significance of Lemme (1.3) is that all obstructions vanish, i.e., d\'eformations of X over base sch\'emas Spec(A/U) (A =local Artin ring with residue corps k) can always be embedded in d\'eformations over Spec(A). Moreover, the theory says that there is a canonical one-one correspondence between Extl(f$\Omega$\_X/k, g)x) and the first order d\'eformations of X, i.e., propre, plat morphismes p, and isomorphismees as follows

\displaycrop[0.72\linewidth]{p079_first_order_deformation.png}

 : X\textasciitilde{} D Xl$\times$ k <\textasciitilde{} X (x P Speck[\textasciitilde{}]/\textasciitilde{}2 D Speck = Speck. Puisque the obstructions vanish, there is a versal formal deJbrmation \textasciitilde{} of X over the base sch\'ema \textasciitilde{}'= Spec Ok[[tl, ..., /hi], where ok=k if the char. is o, or the complete regular local ring, max. ideal p. ok, residue corps k, if char. (k) =p, unique (by Cohen's structure theorem) and N = dimkExtl(g$\Omega$\_X, \textasciitilde{})c)-\end{quote}
8o P. DELIGNE AND D. MUMFORD This means that :Y is a formal sch\'ema, propre and plat over M/, with fibre X over Spec(k) and the two propreties: a) Every deformation of X is induced from YC, i.e., if A is a local Artin 0k-algebra with residue corps k, and p : Y \textasciitilde{} Spec(A) is propre and plat with fibre X over Spec(k), then there is a commutative diagram:

\displaycrop[0.65\linewidth]{p080_versal_deformation_diagram.png}

 Y =\textasciitilde{}/Txj\textasciitilde{}Spec(A) > Spec(A) ] > \textasciitilde{}/ Spec(k) b) If A = k[e]/(\textasciitilde{}2), the above morphismefis uniquely determined by the diagram. This implies that the tangent space to dt'$\times$ at its point ferm\'e is canonically isomorphe to Extl(Ox/k, O\_X). In case Ext\textasciitilde{} g)x)=(o), \textasciitilde{}/\textasciitilde{}r is, in fact, universal: i.e., in proprety a), f is always unique, which means that the foncteur represented by \textasciitilde{} in the cat\'egorie of artin, local ok-algebras is isomorphe to the foncteur of d\'eformations Y/A of X. This fortunately holds for stable curves:

\begin{quote}Lemme (I.t). --- If X is a stable curve, Ext\textasciitilde{} Ox)=(o ).\end{quote}
\begin{quote}D\'emonstration. --- We may assume that k is alg\'ebriquement clos. Now a homomorphismee from f$\Omega$\_X to d)x is given by an everywhere regular vector corps D on X. Such a vector corps is given, in turn, by a regular vector corps D' on the normalisation X' of X which vanishes at all points of X' lying over the points doubles of X. In particular, D' and hence D vanishes identically on all composantes E of X whose normalisation E' has genus >2. There remain the following possibilities for E:

\displaycrop[0.9\linewidth]{p080_component_cases.png}

 iiill E non-singular rational E' rational, E one double pt. E non-singular elliptic E' rational, E>2 double pts. / E' elliptic, E> I points doubles\end{quote}
In all cases where E' is rational, note that D' has to have at least 3 zeroes; and where E' is elliptic, D' has to have at least one zero. So D' vanishes on all composantes E'. This proves that Ext\textasciitilde{} \textasciitilde{})x)= (o). Q.E.D. Schlessinger's theory also allows us to trace what happens to the singularities of X in this deformation Jt. For each point ferm\'e xeX, he studies d\'eformations of the complete local ring gx,x alone, i.e., plat A-algebras 0 plus isomorphismees: 0| $\Omega$\_X,x.

This is a foncteur of A exactly as before. Puisque Ext\&(gl\&/k, g,)=(o), there are no obstructions to extending d\'eformations. First order d\'eformations with A=k[r162 2)

are classified by Ext\&(fl\&i,, 0,). And there is a versal formal deformation \textasciitilde{}, which is a complete local ring and a plat 0k[[tl, ..., tN]]-algebra, N =dimkExt\&(t)\&/k, \textasciitilde{},), such that g/(p, h, . . . , tN)=g.,x. Whenever X is lisse over k at x, N = o, and the theory is uninteresting. The first non-trivial example is a k-rational ordinary double point with k-rational tangent lines: =-kr [u, 073/(,. v). Then N=I, and r v, td]l( V-tl). In other words, if 0 is any deformation of $\Omega$\_X,x and u, v are lifted suitably into g), then u. v= weA and 0 is induced from \textasciitilde{} via the homomorphismee ok[[tl]]$\to$A taking ti to w. This is easy to prove. Finally, Schlessinger's theory connects global d\'eformations to local ones. Soit o\textasciitilde{}/\textasciitilde{},g\textasciitilde{} with \textasciitilde{}'g\textasciitilde{} = Spec(A) be the versal global deformation, let xl, ..., xk\textasciitilde{}X be the points where X is not lisse over k, and let \textasciitilde{}i as Ai-algebra be the versal deformation of the local ring \textasciitilde{},i.x" Soit \textasciitilde{}dc,0=Spec(Al| | Ak). Then we may consider the local rings of \&r at xi. There are o$\simeq$homomorphismees $\varphi$i:A/.$\to$A such that Oi,x\textasciitilde{}=Oi| A. Dualizing, we obtain a morphisme O=H Spec($\varphi$i) : dt'gl-$\to$dt'\textasciitilde{}0 i which describes exactly how the various singularities of X behave in the versal deformation f. The final fact that we need is:

\begin{quote}Proposition (T. 5)- --- $\varphi$ : \textasciitilde{}'gz---\textasciitilde{}'z0 isformally lisse, i.e., there are isomorphismees\end{quote}
,\textasciitilde{}'g, g Spec 0k[[tx, ..., tN+M]]

-g,o \textasciitilde{}Spec o,[[tt, ..., tN]] such that \textasciitilde{}*( tl) =ti, I < i < N.

8\textasciitilde{} P. DELIGNE AND D. MUMFORD D\'emonstration. --- In view of the foncteurial significance of Mt'gz and .At'10, this follows if we prove that the natural map: k (.) Ext$\Omega$\_X (\textasciitilde{}X/k' 0X) \_ 1 \textasciicircum{} is surjectif; (,) is the induced map dap on the tangent spaces to ddgt and \textasciitilde{}'\textasciitilde{}0. Puisque f$\Omega$\_X is an invertible Ox-Module outside the xi's, it follows that: k k 1 \textasciicircum{} A i\textasciitilde{}lEXt$\Omega$\_Xi(\textasciitilde{}2zi, 1\textasciitilde{}.)$\simeq$ I-[ Ext,, (f\textasciitilde{} , \textasciitilde{})\textasciitilde{}i) \textasciitilde{}i- --- i=1 $\Omega$\_Xi\textasciitilde{} xi H\textasciicircum{}0 x, Ext\textasciitilde{},(ax, ex)). Par cons\'equent, (,) is surjectif by the spectral sequence used in Lemme (z .3) and the fact that t-P(X, Ext\textasciitilde{} O\_X))=(o ). Q.E.D. In particular, suppose C is a stable curve over an alg\'ebriquement clos ground corps k, and let xl, ..., xk be the points doubles of C. Soit N=dim Extl(f$\Omega$\_X, O\_X). Then C has a universal formal deformation \textasciitilde{}f/Mr where ..s o,[[tl, ..., t\textasciitilde{}]]. Note that since in this case, the invertible sheaf c\textasciitilde{}/\textasciitilde{} is relatively ample, \textasciitilde{} is non seulement a formal sch\'ema over ..It', but also the formal completion of a unique sch\'ema propre and plat over dr', which we will also denote by \textasciitilde{}. cg is clearly a stable curve over .\textasciitilde{}t'. Now each double point xi has one modulus (cf. our example above) so the versal deformation space of the rings O,i.c is o\textasciitilde{}[[ts ..., t\textasciitilde{}]]. By Proposition (I.5) , we may identify ti with t\textasciitilde{}, and we conclude that for suitable ur v\textasciitilde{}: ..., In particular, tr = o is the locus in \textasciitilde{} where " xi remains a double point " The relation between the formal espace de modules dr' of C and the local structure of H\_g at a point x with $\times$ corresponding to some tri-canonical model of C is exactly the same as in the case of non-courbes singuli\`eres ([M1] , chap. 5, w 2). Soit \textasciitilde{} be the completion of the local ring O,,H\_g, and let T = Spec(\textasciitilde{}). Soit x denote the point ferm\'e of T too. The universal family of stable curves ZgcH\_g$\times$ 5g-6 induces a family Z'cT $\times$ whose fibre Z', over x is C. Then there is a unique morphisme f: T\textasciitilde{}' such that Z'\textasciitilde{}$\times$ , with this isomorphismee restricting to the identit\'e on the fibres over x, both of which are C. I claim that via f, T is formally lisse over ,/t', i.e., O,\textasciitilde{}ok[[tl, ..., tN, tN+l, ..., t\textasciitilde{}]]. In fact, by choosing an isomorphismee p\textasciitilde{} \textasciitilde{}| ( .(\textasciitilde{}r Tphg-6X"ft', we obtain a tri-canonical embedding c\textasciitilde{}fipsg-6$\times$162 of c\textasciitilde{}, hence a morphisme s : \textasciitilde{}'---$\to$H\_g such that \textasciitilde{}, with this embedding, is the pull-back of Zg. Then s factors through T and s : \textasciitilde{}($\to$T is a section off. On the other hand, consider the action of PGL(hg---6) on H\_g. Soit S\textasciitilde{} be the stabilizer of the k-valued point x. Then S\textasciitilde{}is fini and r\'eduites. Because if it were not, S, would have a non-trivial tangent space at the origin, i.e., there would soit un k[z]/(z2)-valued point of PGL(hg---6) centered at the identit\'e, which maps the embedded stable curve C c phg- \textasciitilde{} corresponding to x into itself. But this action is given by an everywhere regular derivation on C, and we have

seen that all such vanish. This means that this k [r / (r -valued automorphisme is the identit\'e at all points of C and, since C is connexes and spans psg-6, the automorphisme is the identit\'e everywhere. Thus S, is fini and r\'eduites. It follows that the action of PGL(5g---6 ) on T is formally free, and hence that T is formally a principal fibre bundle over Mr' with group PGL(5g---6 ). Par cons\'equent T is formally lisse over .W as required. Putting this together with what we know about \textasciitilde{}/t', we conclude the following: Soit k soit unny alg\'ebriquement clos corps, let H'g= H\_g $\times$ Spec(0k), Z'g= Zg $\times$ Spec(ok), let xeH'g soit un point ferm\'e, let C cZ'g be the stable curve over x, let xl, ..., xt\textasciitilde{}eC be its points doubles. Then

\begin{quote}Th\'eor\`eme ( I. 6). --- There are isomorphismees\end{quote}
ok[[t,, ..., g\textasciitilde{}.z; \textasciitilde{} ok[[ui, vr tx, ..., ti\textasciitilde{}]]/(u\textasciitilde{}v,---t,).

\begin{quote}Corollaire (x. 7). --- H\_g is lisse over Z. In particular, for all alg\'ebriquement closcorpss k,\end{quote}
H\_g $\times$ Spec(k) is a disjoint union of a fini number of non-singular algebraic varieties over k. Soit 0\_\{xeH\_g l the corresponding stable curve (Zg)\textasciitilde{} is non-singular\}. H\_g--- S = \{xeZg ] the projection 7: : Zg---$\to$H\_g is not lisse at x\}.

\begin{quote}D\'efinition ( x . 8). --- Soit p : X-$\to$Y soit un lisse morphisme offini type, with Y a noetherian\end{quote}
sch\'ema, and let D cX soit un relative Cartier divisor. Then D has normal crossings relative to Y if for all xED, the local equation d= o of D decomposes in the strict completion (1) "\textasciitilde{},x of 0\textasciitilde{}.x as d=dl.., d\textasciitilde{}, where dl, ..., dk are linearly independent in fftx, x/fft\textasciitilde{}, x + my, y.'\textasciitilde{}, x, with y=p(x).

\begin{quote}Corollaire (x.9). --- H\textasciicircum{}0 *, where S\textasciicircum{}* is a divisor with normal crossings relative\end{quote}
to Z. Zg and S are lisse over Z, and the projection p : S-$\to$S\textasciicircum{}* is fini and an isomorphismee at all points where S\textasciicircum{}* is lisse over Z, i.e., S is the normalisation of S\textasciicircum{}*. ProO --- In the notation of Th\'eor\`eme (1.6), S\textasciicircum{}* is defined in g\textasciitilde{},a\} by the local equation t,... tk= o. And \textasciitilde{},Z\textasciitilde{} '\textasciitilde{} Ok[[Ui, Vi, tl, ..., ti---1, ti+l, ..., tN]] 0k[[tl, ..., t,\_,, t,+,, ..., ts]]. Q.E.D. (x) The complete local ring, formally \'etale over \textasciitilde{}, x with residue corps the separable closure of d)x,x/rex, x.

Next we take up the isomorphismees and automorphismees of stable curves. Supposons que p : X$\to$S, q : Y$\to$S are two stable curves:

\begin{quote}D\'efinition (I. 1o). --- Isoms(X , Y) is thefoncteur on (Sch/S) associating to each S-sch\'ema S'\end{quote}
the set of S'-isomorphismees between X$\times$ and Y$\times$ S'. If X=Y, we denote Isoms(X , X) by Auts(X ). Puisque both X and Y have the canonical polarizations CO\_X/s,Oy/s respectively, any isomorphismee f:X$\to$Y must satisfy J*(c\%:/s)\textasciitilde{}Ox/s. Par cons\'equent, by Grothendieck's results on the representability of the H\textasciicircum{}1lbert sch\'ema and related foncteurs [Grl] , we conclude that Isoms(X , Y) is represented by a sch\'ema Isoms(X , Y), quasi-projective over S. Concerning this sch\'ema, we have:

\begin{quote}Th\'eor\`eme (x. II). --- Isoms(X , Y) is fini and unramified over S.\end{quote}
\begin{quote}D\'emonstration. --- To check that Isoms(X , Y) is unramified, we may take S to be the spectrum of an alg\'ebriquement clos corps k, in which case Isoms(X , Y) is either empty or isomorphe to Aut\textasciitilde{}(X). A point ofAutk(X ) with values in k[r z) with image the identit\'e may be identified with a vector corps on X. By Lemme (i.4), stable curves have no non-zero vector corpss. This proves that Isoms(X , Y) is unramified over S, and since it is also of fini type over S, it is quasi-fini over S. It remains to check that Isoms(X , Y) is propre over S. Locally over S, X and Y are the pull-backs of the universal tri-canonically embedded stable curve by some morphismes from S to H\_g, so that it suffices to prove the propreness of Isoms(X , Y) in the " universal " case where S = H\_g $\times$ H\_g, X and Y being the two image inverses of the universal curve on H\_g. In that case, the ouvert ofIsoms(X , Y) corresponding to lisse curves is dense, so that the Th\'eor\`eme follows from the valuative criterion of propreness which holds by:\end{quote}
\begin{quote}Lemme (I, I2). --- Soit X and Y be two stable curves over a anneau de valuation discr\`ete tl. with\end{quote}
alg\'ebriquement clos residue corps. Denote by \textasciitilde{} and s the generic and point ferm\'es of Spec(P,.), and assume that the fibre g\'en\'eriques X\textasciitilde{} and Y\textasciitilde{} of X and Y are lisse. Then any isomorphismee \% between X\textasciitilde{} and Y\textasciitilde{} extends to an isomorphismee \textasciitilde{} between X and Y. (A posteriori, it follows from Th\'eor\`eme (I. I I) that the lemma holds for any valuation ring P,. and without assuming X, or Y, lisse.) D\'emonstration. --- Another way to put the lemma is that if we start with a lisse curve X, of genus g>2 over the corps des fractions K of I(, there is, up to canonical isomorphismee, at most one stable curve X over R with X, as its fibre g\'en\'erique. We shall deduce this from the analogous uniqueness assertion for minimal models ([L] and [S]): given a lisse curve X\textasciitilde{} of genus g> I over K, there is, up to canonical isomorphismee, at most one regular 2-dimensional sch\'ema X, propre and plat over R, with X\textasciitilde{} as its fibre g\'en\'erique, without exceptional curves of the first kind in X s. Soit z denote a generator of the maximal ideal of R and consider the affine plane curve C n over R given by: xy=z'*.

Soit C. denote the sch\'ema obtained by: I) blowing up the maximal ideal at the unique singularity of Cn; 2) blowing up the maximal ideal at the unique singularity of this sch\'ema, and so on [ l im s It ea , to chec th t regular whos fibre sp\'eciale is the same as that of C n except that the Singular point is replaced by a sequence of n---I projective lines as follows:

\displaycrop[0.62\linewidth]{p085_chain_resolution.png}

 .xz\%z 5 Now suppose x is a singular point of the stable curve X over R.. At x, X is formally isomorphe as sch\'ema over R to one of the sch\'emas Cn, so we may blow up X the same way we blew up Cn. If we do this for all singular points of X, we get a regular sch\'ema with fibre g\'en\'erique X\textasciitilde{}. In addition, any composante rationnelle non singuli\`ere of Xs is linked to the other irr\'eductible composantes by at least two points, hence it is not exceptional of first kind. Par cons\'equent X is the minimal model of X\textasciitilde{}. Note finally that C n is a normal sch\'ema, hence so is X; therefore X is the unique normal sch\'ema obtained from X by contracting all non-singular rational composantes of Xs linked to the other irr\'eductible composantes by exactly two points. This proves that X is essentially unique. Q.E.D. Another important fact about the automorphismees of stable curves is:

\begin{quote}Th\'eor\`eme (1.13). --- Soit k soit unn alg\'ebriquement clos corps and X a stable curve over k.\end{quote}
Soit Pie\textasciitilde{} denote the group of invertible sheaves on X of degree o on each component. Then the map (of ordinary groups) :

Autk(X) $\to$ autk(Pxc (X)) is injective. D\'emonstration. --- Soit $\varphi$soit unn automorphisme of X inducing the identit\'e on PIe\textasciitilde{}

\begin{quote}Lemme ( I. x4). --- If X is lisse, then \textasciitilde{} is the identit\'e.\end{quote}
\begin{quote}D\'emonstration. --- If not, by the Lefschetz-Weil fixed point formula, the number n of fixed points of \% counted with their multiplicities, is n = I --- Tr(\% Tl(Pic\textasciitilde{} + i = 9' - - 2g<o which is absurd. Q.E.D.\end{quote}
\begin{quote}Lemme (x. 15). --- If X is irr\'eductible, then \textasciitilde{} is the identit\'e.\end{quote}
\begin{quote}D\'emonstration. --- Soit $\varphi$' be the action of $\varphi$on the normalisation X' of X. Each singular point of X, together with an ordering of its 2 image inverses in X', defines a distinct morphisme from Gm to Pic\textasciitilde{} so that the image inverse S of the singular locus of X is pointwise fixed by $\varphi$'. One has either a) genus(X')>2: then conclude by Lemme (I. 14); b) genus(X')=i, IS1>2, then $\varphi$' is a translation on X' leaving a point fixed, and so $\varphi$' is the identit\'e; c) X' is the projective line, [ S 1>4 and $\varphi$' is a projectivity leaving more than three points fixed, so is the identit\'e. Q.E.D.\end{quote}
cible composantes of is the case, unless X i points. Q.E.D. Soit I" be the following (unoriented) graph: (i) The set of vertices of F is the set Yo of irr\'eductible composantes of X, (ii) the set of edges of P is the set pl of the singular points of X which lie on two distinct irr\'eductible composantes, (iii) an edge xa I'i has for extremities the irr\'eductible composantes on which x lies.

\begin{quote}Lemme (x. i6). --- If \textasciitilde{} induces the identit\'e on F, then \textasciitilde{} is the identit\'e.\end{quote}
\begin{quote}D\'emonstration. --- If X 1 is an irr\'eductible component of X, then 9(Xi)=X1 and 9 leaves fixed the points of intersection of X 1 with the other composantes. In addition, PIe\textasciitilde{} maps onto PIe\textasciitilde{} so that $\varphi$ acts trivially on Pie\textasciitilde{} Either: a) genus (X1)>2 and 91Xi is the identit\'e by Lemme (I.15) ; b) genus (X1)= i, $\varphi$ acts by a translation and leaves a point fixed, so is the identit\'e on X1; c) X 1 is the projective line and $\varphi$ leaves fixed at least three points, so is the identit\'e on X 1. Q.E.D. Lernma (I. I7). --- (i) Any edge in F has distinct extremities. (ii) Any vertex which is the extremity of o, i or 2 edges is fixed by 9. (iii) \textasciitilde{} acts trivially on H\textasciicircum{}0(F, Z). D\'emonstration. --- It is easy to check that the subgroup of PIc\textasciitilde{} corresponding to invertible sheaves whose restriction to each irr\'eductible component of X is trivial is canonically isomorphe to H\textasciicircum{}0(F, Z)| m. This implies (iii), and (i) is trivial. The morphisme from PIe\textasciitilde{} to the product OPIe\textasciitilde{} extended over the irreduX, is surjectif, so that if Pie\textasciitilde{} + \{e\}, then $\varphi$(Xi)=X i. This is a projective line, linked to the other composantes in at least three We prove now that if an automorphisme 9 of any fini graph P has the propreties stated in Lemme (I. I7) , it is the identit\'e. Make induction on the sum of the number of vertices and edges ofF. IfP has an isolated point x, then $\varphi$(x)=x so let F* = F---\{x\}. Then $\varphi$= identit\'e on F* by induction, so 9 = identit\'e on F too. If F has an extremity x, then $\varphi$(x)----x, and again let P* be F minus x and the edge abutting at x. Then F* has all the propreties P has, so $\varphi$= identit\'e on P*, hence $\varphi$=identit\'e on P. If F has a vertex x on which only 2 edges abut, we have one of the two cases:

\displaycrop[0.9\linewidth]{p086_graph_cases.png}
\displaycrop[0.8\linewidth]{p087_fixed_edge_reduction.png}

 (2) \_ \_ \_ e1 e\textasciitilde{}\end{quote}
In the first case, $\varphi$(y)=y, and let F* be F minus x, e1and e2. Then $\varphi$=identit\'e oil F* and $\varphi$(ei)= e\textasciitilde{} too, since if $\varphi$reverses the e\textasciitilde{}'s, this contradicts (iii). In the second case, let I'* be P minus x, and with e1 and ez identified: e Then $\varphi$= identit\'e on P*, so $\varphi$= identit\'e on P. Next, say F has an edge e with extremities x andy such that $\varphi$(x)=x, $\varphi$(y)=y, $\varphi$(e)=e. Soit P* be P minus e. Then $\varphi$=identit\'e on F*, so $\varphi$= identit\'e on P. If none of these reductions are possible, we must be in a situation where a) every vertex is the abutment of at least three edges and b) no edge is left fixed. It is easily seen that the first Betti number b1 of any of the connexes composantes of I' is at least 2. Soit no= number of fixed vertices nl=number of edges reversed by \textasciitilde{}. Then, unless F=O, the Lefschetz fixed point formula reads: no + nl = b0---bl<O , which is impossible. Q.E.D.

\section*{2. D\'eg\'en\'erescences de courbes et de leurs jacobiennes}
We consider the situation: K = discretely-valued corps; 1\textasciitilde{} = integers in K, k = l\textasciitilde{}/\textasciitilde{}Jl = residue corps (assumed alg\'ebriquement clos) ; S =Spec(R), \textasciitilde{} and s its generic and point ferm\'es respectively; C = a curve, lisse, geometrically irr\'eductible and propre over K, of genus g\textasciitilde{} 2 ; J = the jacobian variety of C; og = the N6ron model of J over R (cf. [N]); J\textasciitilde{} the open subgroup sch\'ema with J\textasciitilde{} component of Js; c\textasciitilde{} = the minimal model of C over R.. A word about the existence and uniqueness of c\textasciitilde{} is needed. We recall that c\textasciitilde{} is to soit un regular sch\'ema, plat and propre over R., with fibre g\'en\'erique c\textasciitilde{}n= C such that t for any other regular sch\'ema c\textasciitilde{},, plat over P,., with fibre g\'en\'erique \textasciitilde{}fn = C, the birational map c\textasciitilde{},\textasciitilde{}c\textasciitilde{} is a morphisme, \v{S}afarevich in IS] and Lichtenbaum [L] have proven that such a \textasciitilde{} (which is obviously unique) exists, provided that there is some regular 5', propre and plat over R, with fibre g\'en\'erique C. And, in fact, that c\textasciitilde{} is projective over R. To construct such a c\textasciitilde{},,proceed as follows: first let c\textasciitilde{},, soit unny sch\'ema, projective and plat over R with fibre g\'en\'erique C. Soit R be the completion of R, and let c\textasciitilde{},,\_\_c\textasciitilde{},,$\times$ spooRSPec R. Then O' is an excellent surface, so by [Ab] and by unpu-

blished results of H\textasciicircum{}1ronaka, there is a sheaf of ideals 3 with support in the singular locus of O' such that blowing up J leads to a regular surface O. But then for some n, so .\textasciitilde{} is induced by a unique Sheaf of ideals Jc0\textasciitilde{}e,,. Soit \textasciitilde{}' be obtained by blowing up J. Then \textasciitilde{}, \textasciitilde{} cg, $\times$ spe,$\times$Spec R. so cg, is regular, c\textasciitilde{}, is also projective over R, hence a projective cg exists.

\begin{quote}D\'efinition (2. x ). --- J has stable reduction if J8 has no radical unipotent.\end{quote}
\begin{quote}D\'efinition (2.2). --- C has stable reduction in sense i if c\textasciitilde{}s is r\'eduites and has only\end{quote}
points doubles ordinaires. C has stable reduction in sense 2 if there is a stable curve c\textasciitilde{}, over R with fibre g\'en\'erique cgs Note that if a stable \textasciitilde{} exists, then by Th\'eor\`eme (I. I I) it is unique.

\begin{quote}Proposition (2.3). --- The two senses of stable reduction for C are equivalent.\end{quote}
\begin{quote}D\'emonstration. --- Say a stable cg, exists. Blowing up the singularities of c\textasciitilde{}, as in\end{quote}
\begin{quote}Lemme (i. 12), we obtain the minimal model \textasciitilde{} of C and it is seen that c\textasciitilde{}8 is r\'eduites\end{quote}
with only points doubles ordinaires. Conversely, suppose the minimal model \textasciitilde{} has this proprety. Soit El, ..., E, be the non-singular rational composantes of \textasciitilde{} which meet the other composantes in only two points. Then the Ei's divide into several chains of the type:

\displaycrop[0.75\linewidth]{p088_component_chain.png}

 unless the entire fibre consisted of E\textasciitilde{}'s and has the type: n>2 \textasciitilde{}'8 = loop of Ei's (E$\varphi$)oo =---2. But in this case genus(C) = genus(Cgs) = I, which contradicts our assumption. Now, according to Th\'eor\`eme (27. I) of Lipman [Li] (generalizing a result of Artin [A2], which works for surfaces of fini type over a corps) any set ofk non-singular rational curves connexes in a chain as above on a regular surface X, with self-intersection 2 on X, can be blown down to a rational double point P of type Akon a normal surface X 0. Reversing the process and blowing up a rational double point of type Ak, it is easy to see that nonsingular branches y on X, crossing transversally only the first or the last rational curve in the chain: E2 Ek . . . .

are still non-singular branches when mapped to X0; and if Y1, Y2 intersect E1 or Ek in distinct points, then they cross transversally on X o. Par cons\'equent, suppose we blow down all the chains of Ei's on c\textasciitilde{}. Soit c\textasciitilde{}, be the normal surface so obtained. Then if one of these chains fits into \textasciitilde{}, like this:

\displaycrop[0.75\linewidth]{p089_contraction_chain.png}

 then on \textasciitilde{};, the images of F and G still have only points doubles ordinaires, each has one non-singular branch through the singular point P, and these branches cross transversally. Par cons\'equent \textasciitilde{}' is a stable curve over R. Q.E.D. We are now ready to prove the key result on which our proof of irreducibility depends:

\begin{quote}Th\'eor\`eme (2.4). --- J has stable reduction if and only if C has stable reduction.\end{quote}
\begin{quote}D\'emonstration. --- The connection between \textasciitilde{} and of is based on the following result of Raynaud [R] :\end{quote}
\begin{quote}Th\'eor\`eme (2.5). --- If c\textasciitilde{} and of are as above, and the greatest common denominator d of\end{quote}
the multiplicities of the composantes of \textasciitilde{} is i, then ofo represents thefoncteur Pic\textasciitilde{} (This result is not stated as such in JR]. It comes out like this, in the terminology o f is paper : a) Condition (N) is verified and p,(0\textasciitilde{}e)-=0s, so b) \textasciitilde{} is cohomologically plat over S in dim. o by Th\'eor\`eme 4; c) therefore Pie\textasciitilde{} is representable and s\'epar\'e over S by Th\'eor\`eme 3; d) since E=\{o\} and Q=P/E, we find P\textasciitilde{} and since C\textasciitilde{}is 1-dimensional over S, p0 and Q0 are lisse over S; therefore R\textasciitilde{} Q0 and P,,= R. e) Then by Th\'eor\`eme 5, Pie\textasciitilde{} is the identit\'e component of the Ndron model of Pic\textasciitilde{} =J.) Now assume that C has stable reduction. Then \textasciitilde{} is r\'eduites so d=I. By

\begin{quote}Th\'eor\`eme (2.5) , of\textasciitilde{}176 Par cons\'equent of\textasciitilde{}176 ). Puisque \textasciitilde{}8 is r\'eduites\end{quote}
with only points doubles ordinaires, its generalized jacobian Pic\textasciitilde{} is an extension of an abelian variety by a torus, i.e., has no radical unipotent. Par cons\'equent J has stable reduction. The converse is more difficult. Assume J has stable reduction. We first prove that C has stable reduction under the additional hypothesis that C has a K-rational point (1). In this case, \textasciitilde{} has an R-rational point, and since \textasciitilde{} is regular, sections of over R pass through composantes of \textasciitilde{} of multiplicity one. Par cons\'equent d=I, and (1) This is in fact the only case which will be needed in our application to questions of irreducibility.

9 \textasciitilde{} P. DELIGNE AND D. MUMFORD

\begin{quote}Th\'eor\`eme (2.5) applies. In particular, Pie\textasciitilde{} jo so Pie\textasciitilde{} ) has no unipotent\end{quote}
radical. We apply:

\begin{quote}Lemme (2.6). --- Soit D soit un complete i-dimensional sch\'ema over k such that H\textasciicircum{}0 \_k,\end{quote}
and such that the generalized jacobian of D has no radical unipotent. Then (i) \%(OD)=X(ODred); (ii) the singularities of Dr0a are all transversal crossings of a set of non-singular branches (i.e., analytically isomorphe to the union of the coordinate axes in Am). D\'emonstration. --- Soit J c OD be the ideal of nilpotent elements. Filtering J by a chain of ideals ork such that J.JkCJk+l, and using the exact sequence: o ---, *k/Jk+lo (r ---" 0 it is easy pour d\'eduire that Pie\textasciitilde{} is an extension of Pie\textasciitilde{} by a unipotent group. Par cons\'equent, since by assumption Pie\textasciitilde{} has no unipotent subgroups, Pic\textasciitilde{} \textasciitilde{}=Pic\textasciitilde{} Puisque H\textasciicircum{}0(OD), resp. Ht(OD\textasciitilde{}ea), is naturally isomorphe to the Zariski tangent space to Pic\textasciitilde{} resp. Pie\textasciitilde{} it follows that H\textasciicircum{}0(0D) \textasciitilde{} Hl(\textasciitilde{})Dred) hence (i) is proven. Soit \textasciitilde{} : C---\textasciitilde{}D$\pi$a be the normalisation of Dr,d and let D* be the local ringed space which, as topological space is D, and whose structure sheaf is given by: F(U, OD.)=\{f\textasciitilde{}F(U, 7L(Oc)) ]x\textasciitilde{}, x2$\simeq$\textasciitilde{}(U),f(x\textasciitilde{})=f(x2) if r:(x\textasciitilde{})=\textasciitilde{}z(x2) \} It is easy to check that D* is a 1-dimensional sch\'ema whose singularities are all transversal crossings of a set of non-singular branches and that \textasciitilde{} factors: C ---$\to$ ---$\to$Dred. I claim that rd' : D*$\to$D$\pi$d is an isomorphismee. Filter OD,/OD\textasciitilde{}eGSO as to obtain a chain of coherent OProd-algebras: (\textasciitilde{}D* : ('0(0) :\textasciitilde{} \textasciitilde{})(1)2\textasciitilde{} . . . \textasciitilde{}1 \textasciitilde{}(k)\textasciitilde{}\_\_. \textasciitilde{})Dred such that l((r162 '\textasciitilde{}+1))=I. Equivalently, this factors rd': D*= DO$\to$ D1---$\to$... \textasciitilde{} Dk= D\textasciitilde{}ea where O(h)\textasciitilde{}OD,, and all arrows are homeomorphismes. If \{x\textasciitilde{}\}=Supp(O(')/O(\textasciitilde{}+\textasciitilde{}l), then we get an exact sequence: O-$\to$ \textasciitilde{})Dn+l-$\to$ ODn-'+kzn$\to$O (kv denotes the residue corps aty, as sheaf on D), and hence: 1 * o \textasciitilde{} k $\to$ n (OD,+\textasciitilde{})$\to$ H\textasciicircum{}0(\textasciitilde{})]\textasciitilde{},)$\to$ o. It follows easily that Pie\textasciitilde{} is an extension of Pie\textasciitilde{} by G\textasciitilde{}. But Pie\textasciitilde{} has no unipotent subgroups, and this can only happen if D\textasciitilde{},a= D*. Q.E.D. for lemma.

We apply the lemma to c\textasciitilde{}. According to Lichtenbaum [L] and \v{S}afarevich [S], there is a divisor K on \textasciitilde{} such that for all positive divisors D lying over the point ferm\'e of Spec(R), we have: (D. (D + kK)) z(eD) -

Soit El, ..., E,, be the composantes of c\textasciitilde{}, dl' ..., dn their multiplicities. Then conclusion (i) of the lemma implies that: But ((\textasciitilde{}diEi).Ek)=o , all k, since \textasciitilde{}diE i is the divisor of a function r:\textasciitilde{}R if (=)----=maximal ideal of R. Par cons\'equent : (*) ((\textasciitilde{}(d,---I).Ei).K)=(\textasciitilde{}Ei.\textasciitilde{}Ei). Note that at least one di equals i since c\textasciitilde{} has a section over Spec(R) and every section must pass through a component of c\textasciitilde{}s of multiplicity I. Moreover the intersection matrix (Ei. ES) is negative indefini, with one-dimensional degenerate subspace generated by \textasciitilde{}]d\textasciitilde{}Ei, hence \textasciitilde{} some di>I , it follows that (\textasciitilde{}E,.\textasciitilde{}Ei)<o. Par cons\'equent, by (*), \$ (Ei0.K)<o for some i0. Then we have: a) (E\textasciitilde{}o.K) <o; b) (Ei0.E\textasciitilde{}o ) < o; c) hence in fact (E\textasciitilde{}0.E\textasciitilde{}0) = (E\textasciitilde{}0.K) =--- \textasciitilde{} so E\textasciitilde{}o is an exceptional curve of the first kind. This contradicts our assumption that on \textasciitilde{} all possible curves have been blown down. Par cons\'equent di= i, all i. This proves that \textasciitilde{} is r\'eduites. By conclusion (ii) of the lemma, plus the fact that the dimension of its Zariski tangent-space is everywhere one or two (since c\textasciitilde{} lies on a regular surface \textasciitilde{}), we deduce that c\textasciitilde{}s has only points doubles. This proves that C has stable reduction in the case that C(K)4=0. In the general case, C will acquire a rational point in a fini extension K' of K. Soit S' be the spectrum of the localisation at some maximal ideal of the cl\^oture int\'egrale of R in K'; S' is the spectrum of a valuation ring and is faithfully plat over S. We put S"=S

\displaycrop[0.62\linewidth]{p091_descent_base_change.png}

'$\times$ and denote by C' and C" the image inverses of C on ! S\textasciitilde{} = Spec(K') t! and S\textasciitilde{} respectively. C" \textasciitilde{} C' prl S" \textasciitilde{} S' p$\pi$ :, C

io > S

Soit C' be the stable curve on S' having C' as fibre g\'en\'erique. The restriction, C', of C' ! to S\textasciitilde{} carries a descent datum with respect to p, i.e. an isomorphismee, \%, between the restrictions of p$\pi$(C') and p$\pi$(C') to S'\textasciitilde{}'. It remains only to extend the isomorphismee $\varphi$ to an isomorphismee, p, between the p$\pi$(C'). Then $\varphi$ will soit un descent datum for C' with respect to p. As C' is canonically polarized (I.2), this descent datum will be effective, and so define a stable curve, C, over S with fibre g\'en\'erique C. Because j0 has no radical unipotent, the image inverse, p,j0, of j0 on S' is the identit\'e component of the N6ron Model of the jacobian of C', so that, defining q=poprl=popr2, one has PicO(C,/S ,) \_p*jO Pic0(p$\pi$,)= 9 0 *---, Pie (pr2C)= q.j0 We denote by T the closed subsch\'ema of Isom(S", p$\pi$(C'), p$\pi$(C')) corresponding to those isomorphismees which induce (via the preceeding identifications) the identit\'e on the image inverse of j0. By (I. I I ), T is fini and unramified over S", and by (I. 13) T is radicial over S". We conclude that the morphisme from T to S" identifies T with a closed subsch\'ema X /! tt of S". As X contains S,, which is schematically dense in S , we have that X is S" and T " is " the desired section, 9, of Isom(S", *---' p$\pi$(C ), pr;(C')) over S" which extends \textasciitilde{}\%. Q.E.D. Combining Proposition (2.3), Th\'eor\`eme (2.4), and the th\'eor\`eme de r\'eduction stable for vari\'et\'es ab\'eliennes quoted in the introduction, we obtain the most important consequence:

\begin{quote}Corollaire (2.7). --- Soit R soit un anneau de valuation discr\`ete with quotient fMd K. Soit C soit un\end{quote}
lisse geometrically irr\'eductible curve over K of genus g > 2. Then there exists a extension alg\'ebrique finie L of K and a stable curve WL over RL, the cl\^oture int\'egrale of R in L, with genericfibre CCl\textasciitilde{},\textasciitilde{} \textasciitilde{}C XKL. § 3" Elementary derivation of the theorem. Soit k soit unn alg\'ebriquement clos corps of char. p 4:o. We use the notation of w I, except that we will now denote by H\_g the product H\_g$\times$ of the previous H\_g with Spec(k): it is a disjoint union of non-singular varieties H\_g,1, ..., H\_g,n over k and 9 Pg 0 is the subsch\'ema of H\textasciicircum{}1lbp0g-0/k of tri-canonical stable curves. Similarly, H a is the open dense subset of H\_g of tri-canonical non-courbes singuli\`eres. By the results of [M1] , we know that a coarse geometric quotient o\_ HO/PGL (5g\_ 6) M\_g--- exists, that it is a disjoint union of normal varieties over k and is the coarse espace de modules * 0 for non-courbes singuli\`eres of genus g. Soit S = H\_g---H\_g. Then everything decomposes into the same set of composantes.

Soit H\_g,i= composantes of H\_g 0 \_ H \textasciitilde{} of H\_g then H\_g,i--- H\_g, i ca 0 and 0 \_ 0 M\_g,i--- H\_g, dPGL(5g--- 6) ------composantes of M \textasciitilde{} and S\textasciicircum{}*------disjoint union of \$1, ..., Sn, Si---\textasciitilde{}H\_g, icaS\textasciitilde{}. We want to prove that M\_g or equivalently H \textasciitilde{} or equivalently H\_g is irr\'eductible. We shall use: (i) the fact that these statements are true in char. 0; (ii) the inductive assumption that these statements are true for smaller genus. Step I. --- No component of M \textasciitilde{} is complete (i.e., propre over k). D\'emonstration. --- Here we use the char. 0 result. By [M1] , there is a sch\'ema X, quasiprojective over Spec(W(k)), W(k) the Witt vectors, whose fibre ferm\'ee is M \textasciitilde{} and whose fibre g\'en\'erique X\textasciitilde{} is the char. 0 coarse espace de modules over the corps des fractions of W(k). In particular, X, is known to be connexes. Puisque X is quasi-projective over W(k), we can embed X as an open dense subset of a sch\'ema X projective over W(k). X\textasciitilde{} is still connexes, hence by the th\'eor\`eme de connexit\'e of Enriques-Zariski [EGA 3], the fibre ferm\'ee "\textasciitilde{}0 of .\textasciitilde{} is connexes. But if Y were a complete variety which is a component of M\_g then: a) Y is an ouvert of M\_g, which is an ouvert of "\textasciitilde{}0, and: b) since Y is propre/k, Y would soit un closed subset of "\textasciitilde{}0 too. Par cons\'equent, Y = "X0, hence M\_g is itself irr\'eductible and complete. On the other hand, if Ag is the coarse espace de modules of principally pclarized g-dimensional vari\'et\'es ab\'eliennes, then the map associating to each curve its jacobian defines a morphisme: 0 : M\_g If M\_g were complete, the image of 0 would be closed. But it is well known that the closure of the image of 0 contains all products of lower dimensional jacobians too, so it is not closed. Q.E.D. Step 1I. --- No component of H\_g consists entirely of non-courbes singuli\`eres, i.e., S\textasciitilde{}\textasciitilde{}=O for all i. D\'emonstration. --- Here we combine Step I with the result of § 2. Take any i. Soit T = Speck [It]]. Puisque M \textasciitilde{}\textasciitilde{}is not complete, there is a morphisme $\varphi$of the generic point T, of T into M\_g which does not extend to a morphisme of T into M\_g Now we replace T t 0 by its normalisation T' in a extension alg\'ebrique finie and let $\varphi$' : T$\simeq$-$\to$M\_g.i be the

induced morphisme, which still does not extend to a morphisme from T' to M\_g.\textasciitilde{}. By the results of w 2, if T' is chosen suitably there exists a stable curve 7: : C'-$\to$T' over T' t whose fibre g\'en\'erique C, is a non-singular curve corresponding to the morphisme r t 0 $\varphi$ : T,---\textasciitilde{}M\_g via the foncteurial propreties of the espace de modules. Puisque T' is the spectrum of a local ring, we can choose an isomorphismee P ($\pi$.(co\textasciitilde{}/3T,))\textasciitilde{}P\textasciitilde{}g-G$\times$ and get a tri-canonical embedding C'cPSg-\textasciitilde{}$\times$ '. C', with this embedding, is then

induced from the universal tri-canonically embedded stable curve by a morphisme t : T'\textasciitilde{}H\_g. Puisque the fibre g\'en\'erique of C' is C,, we get a commutative diagram:

\displaycrop[0.72\linewidth]{p093_lift_to_hilbert_diagram.png}

 + T' ) H\_g U U $\pi$s",Ol T', \_ 0 \textasciitilde{} H\_g M \textasciitilde{} C: o " 9 \textasciitilde{} M9 If Im(qJ) cH \textasciitilde{} then $\varphi$' would extend to a morphisme from T' to M \textasciitilde{} hence from T' to M\_g and this is a contradiction. Moreover, \textasciitilde{}(T,) must soit un point of H\textasciicircum{}0 since its image in M \textasciitilde{}is in M \textasciitilde{} Par cons\'equent the image x of the point ferm\'e of T' by \textasciitilde{} is in the closure

of H\_g,i, i.e., in H\_g, i, but not in HOg./itself. Q.E.D. Step IIL --- S\textasciicircum{}* is connexes. D\'emonstration. --- This will follow using only the induction assumption of irreducibility for lower genera. Soit Z cH\_g$\times$ 5g-G be the universal tri-canonically embedded stable curve. Soit ScZ be the set of points where Z is not lisse over H\_g. As we proved in w I, S is non-singular and is the normalisation of S\textasciicircum{}*. In particular, this shows that if xeS\textasciicircum{}*, then the corresponding curve Zx has exactly one double point if and only if x is a non-singular point of S\textasciicircum{}*. Stable curves C of genus g with exactly one double point belong to one of the following types:

\displaycrop[0.5\linewidth]{p094_curve_types.png}

 type o : C\textasciitilde{} type k: C irr\'eductible normalisation C' of C has genus g---I. C has two non-singular composantes C1, C2 genus(C1) =k genus(Ca) =g--- k ra] if one dou 'e oint is o ty o ' hon open dense subset of S\textasciicircum{}* of non-singular points is the disjoint union of ouverts S\textasciicircum{}*(o),...,S\textasciicircum{}*(/g/). We first check: (.) Each set S\textasciicircum{}*(k) is irr\'eductible. D\'emonstration of (.).- Take the case k=o: the cases i<k< [g] are similar. Soit T= \{(xl, x2) ix 1+x2 and $\pi$(x,)=$\pi$(x\textasciitilde{})eH\textasciicircum{}0\_,\}cZg\_ 1XH\_g\_lZg\_l. T is lisse with irr\'eductible fibres over H\textasciicircum{}0 hence T is irr\'eductible. Consider the correspondence relating S\textasciicircum{}*(o) and T:

,'set of pairs

\displaycrop[0.9\linewidth]{p094_correspondence_diagram.png}

 xeS\textasciicircum{}*(o), \{xl, x2\}eT such that if y = re(x,) = =(x2), j I then there exists a birational morphisme \textasciitilde{}such that f(xl)=f(x2). It is easy to check that W is Zariski-closed. Moreover, for any \{x,, x2\}eT , W n (S\textasciicircum{}*(o) $\times$ \{xl, x2\}) is an orbit in S\textasciicircum{}*(o) under PGL(5g---6): so these are non-empty irr\'eductible subsets all of the same dimension. It follows that W itself is irr\'eductible. But the projection from W to S\textasciicircum{}*(o) is surjectif, so S\textasciicircum{}*(o) is irr\'eductible. Q.E.D. for (,). Now for any k, I<k< lg/, choose any stable curve C(k)of the type:

\displaycrop[0.85\linewidth]{p095_stable_curve_Ck.png}

 O Q, one non-singular component C(k)' of genus k. Q2 one component C(k)" with one double point, normalisation of genus g---k---I. Soit P(k) soit un point of H\_g such that Zp(k)-\textasciitilde{}C(k). Step III will be completed if we prove : (**) P(k) is in the closure of S\textasciicircum{}*(o) and of S\textasciicircum{}*(k), hence S\textasciicircum{}*(o), S\textasciicircum{}*(k) both lie in the same topological component of S\textasciicircum{}*. D\'emonstrationof (**). --- Soit T= Spec k[[t]]. Using the fact that S\textasciicircum{}* has two branches through P(k), one for each of the points doubles of C(k), we see that there exist two morphismes fl,f2

\displaycrop[0.68\linewidth]{p095_degeneration_pair.png}

 : T$\to$S\textasciicircum{}* f,(T\textasciitilde{}) =f\textasciitilde{}(T\textasciitilde{}) = P(k), T\textasciitilde{} = point ferm\'e of T D\textasciitilde{},,q D2, ,,1 ;, / ,/ / ,.J--- ,, <--- , C(k) C(k) such that if =1: DI$\to$T, =2:D2$\to$T are the two stable curves over T induced by fl and f2, then: a) the fibre ferm\'ees Dl, s, Dz, 8 are C(k) ; b) there are sections sx : T$\to$D1, s2:T$\to$D 2 whose images are non-lisse points of re1,\textasciitilde{} and such that sl(Ts) and s\textasciitilde{}(T\textasciitilde{}) are the two points doubles Q.1 and Qa of C(k) respectively, and : c) the fibre g\'en\'eriques DI,,, D2,\textasciitilde{} have only one double point (el. figure). I claim: A) D1,, is of type (k); B) D2,\textasciitilde{} is of type (0).

To prove A), let D'i be the result of blowing up the subsch\'ema s1(T) of D1. Then D[ is still plat and propre over T and its fibre sp\'eciale is the fibre sp\'eciale of D1with Q1 blown up (use the fact that formally at Qt, D1 is isomorphe to k[[t, x,y]]/(x.y) with the section given by x=y=o). Par cons\'equent the fibre sp\'eciale of D'1is the disjoint union of C(k)', C(k)", so the general fibre of D'1 is the disjoint union of two irr\'eductible curves which specialize to C(k)', C(k)" respectively. Puisque D\textasciitilde{},\textasciitilde{} has only one double point, D'i,\textasciitilde{} is non-singular, so D'I,\textasciitilde{} is the disjoint union of two non-singular irr\'eductible curves which must then have the same genera as C(k)' and C(k)", i.e., k, g---k. Thus D1,\textasciitilde{} has type (k). To prove B), it suffices to check that D2,\textasciitilde{} is geometrically irr\'eductible. If not, D2,n would have two composantes meeting at the single point sz(Tn). Puisque D2 is lisse over T at each generic point of its fibre sp\'eciale, distinct geometric composantes of D2, have to have specializations which are distinct composantes of (Dz),. Then (D2)8would have two composantes meeting at the point Qz = s2(Ts). This is false, so B) is proven. Now because of A), fl(T\textasciitilde{})eS\textasciicircum{}*(k), hence P(k)=fl(T,) is in the closure of S\textasciicircum{}*(k). And because of B), f2(T\textasciitilde{})cS\textasciicircum{}*(o), hence P(k)=f2(T\textasciitilde{}) is in the closure of S\textasciicircum{}*(o) too. Q.E.D. for (**). This completes the proof of Step III since we now see that all irr\'eductible composantes of S\textasciicircum{}* are part of the same topological component. Finally, from Steps II and III, we see that a) S\textasciicircum{}*, being connexes, is part of a single component of H\_g, while b) each component of H\_g contains part of S\textasciicircum{}*. Thus H\_g is irr\'eductible, as was to be proven. § 4"

\displaycrop{p098_morphism_stack_square.png}
\displaycrop{p098_stack_diagonal_display.png}
\displaycrop[0.45\linewidth]{p100_base_change_square.png}
\displaycrop[0.72\linewidth]{p100_property_diagram.png}
\displaycrop[0.45\linewidth]{p100_morphism_square.png}
\displaycrop{p101_chow_lemma_diagram.png}
\displaycrop{p101_more_display_diagram.png}
\displaycrop[0.45\linewidth]{p103_composition_diagram.png}
\displaycrop[0.55\linewidth]{p103_stack_atlas_diagram.png}
\displaycrop[0.65\linewidth]{p103_long_diagram.png}
\displaycrop[0.6\linewidth]{p105_base_change_diagram.png}

 Some results on champs alg\'ebriques. The proofs of the results stated dans cette section seront donn\'es ailleurs. Soit C soit un cat\'egorie and let p : 5\textasciitilde{}C soit un cat\'egorie over C. For each UeOb C, we denote by 5\% the fibre p-l(U). The cat\'egorie 5P is fibered in groupoids over C if the following two conditions are verified: a) For all 9:U$\to$V in C and yeObS\textasciitilde{} there is a map f:x$\to$y in 5P with p(f) = \textasciitilde{}. b) Given a diagram in \textasciitilde{}9\textasciitilde{} let x

y/ U / V

be its image in C. Then for all x:U\textasciitilde{}V such that $\varphi$=+X, there is a unique h :x$\to$y such that f=g.h and p(h)=x. Condition b) implies that the f: x$\to$y whose existence is asserted in a) is unique up to canonical isomorphismee. Assume that for each $\varphi$ : U$\to$V in C and each YeOb 5:v, such an f: x\textasciitilde{}y has been chosen. This x will be written as $\varphi$*y. Then, $\varphi$*" is " a foncteur from 5:v to 5: U and if $\varphi$+ is a composite morphisme in C, the foncteurs ($\varphi$+)* and +*$\varphi$* are canonically isomorphe. We propose the terminology " champ " for the French word " champ " of nonabelian cohomology (Giraud [G]).

\begin{quote}D\'efinition (4. I ). --- Soit C soit un cat\'egorie with a Grothendieck topology. We assume products\end{quote}
and fibre products exist in C. A champ in groupoids over C is a cat\'egorie over C,p : 5r such that: (i) 5: is fibered in groupoids over C. (ii) For any UeOb C and any objets x, y in 5\textasciitilde{}j the foncteur from C/U to (sets) which to any $\varphi$:V\textasciitilde{}U associates H\textasciicircum{}0ms:v($\varphi$*x, $\varphi$*y) is a sheaf. (iii) If $\varphi$\textasciitilde{}:Vi\textasciitilde{}U is a covering family in C, any descent datum relative to the $\varphi$i, for objets in 5", is effective. For each xaOb \textasciitilde{}9\textasciitilde{} there are given isomorphismees between the image inverses of xi= $\varphi$ix and xj=$\Omega$\_X over Vo=Vi$\times$ and the pull-backs of these isomorphismees on Vi/k=Vi$\times$ \textasciitilde{} satisfy a " $\omega$\_cycle " condition. In (iii) it is required that reciprocally, any such " descent datum " be defined by some xaS\textasciitilde{} . In whatfollows, for the sake of brevity, we will use " champ " to mean " champ in groupoids " If U cOb C and if \textasciitilde{}9 \textasciitilde{} is a champ over C, the fibre 5:U will be called the cat\'egorie of sections of 5e over U. Soit C soit uns in (4. I)- The champs over C are the objets of a 2-cat\'egorie [B] (champs/C): i-morphismes are foncteurs from one champ to another, compatible with the projection into C; and 2-morphismes are morphismes of foncteurs. In this 2-cat\'egorie every 2-morphisme is an isomorphismee. Products and 2-fibre products exist in this 2-cat\'egorie. To each XeOb C is associated the " representable " champ over C whose cat\'egorie of sections over U is the discrete cat\'egorie whose objets are the morphismes from U to X. This champ will be denoted simply X. For any champ 5e, the cat\'egorie H\textasciicircum{}0m(X, 5:) is canonically equivalent to the cat\'egorie of sections of 5\# over X. Because of this, 5e is sometimes said to " classify " its sections over variable XeOb C. In the cat\'egorie H\textasciicircum{}0m(5", X), all morphismes are identities, i.e., H\textasciicircum{}0rn(5:, X) is just a set. Soit us denote also by C the 2-cat\'egorie having the same objets and morphismes as C, and in which the identities are the only 2-morphismes. The above construction then identifies C with a full sub-2-cat\'egorie of (champs/C).

For each S\textasciitilde{}Ob C, the cat\'egorie C/S satisfies the assumptions of (4.I). Any champ 5o0 over C/S (an " S-champ ") gives rise to a champ 5\textasciitilde{} over C; a section (\% \textasciitilde{}) of 5O over U\textasciitilde{}ObC consists of (i) a morphisme $\varphi$ : U\textasciitilde{}S; (ii) a section \textasciitilde{} of 5o0 over (U, 9)-

\begin{quote}D\'efinition (,t.2). --- A I-morphisme of champs over C, F : 5O1\textasciitilde{}5O2, will be called repre-\end{quote}
sentable if for any X in C and any i-morphisme x : X$\to$5O2, thefibre product X$\times$ is a representable champ. In down to earth terms, this means the following: (i) for any f: Y$\to$X in C, the cat\'egorie whose objets are pairs \{a section, y, of 5\textasciitilde{} over Y; an isomorphismee F(y)-\textasciitilde{}f*(x)\} is equivalent to a cat\'egorie S(f) in which all morphismes are identities; (ii) the foncteur f\textasciitilde{}Ob S(f) is representable by some g : Z$\to$X. Such a Z represents the fibre product X $\times$162 Soit/' soit un proprety of morphismes in C, stable by change of base and of a local nature on the target.

\begin{quote}D\'efinition (4.3). --- A representable morphisme F : 5O1$\to$5O2 of champs over C has proprety t"\end{quote}
if for any I-morphisme x : X$\to$SO 2 the morphisme in C deduced by changement de base: F' : X$\times$ SI$\to$X has that proprety.

\begin{quote}Proposition (4.4). --- Soit 5O soit un champ. The diagonal map\end{quote}
5O-'. SOx so is representable if and only if for all X, Y\textasciitilde{}Ob C and i-morphismes x : X$\to$SO, y : Y$\to$SO, the fibre product X x\textasciitilde{}Y is representable. If X\textasciitilde{}Ob C and x, y are sections ofsoover X, we denote by Isom(X, x,y) the sheaf on C/X which to every Z over X associates the set of isomorphismees between the image inverses ofx andy over Z. Then the objet representing 5O$\times$215 (the prodtlct taken with the map (x,y) :X\textasciitilde{}SP$\times$ is just Isom(X,x,y). If x:X\textasciitilde{}so, y :Y\textasciitilde{}SO are I-morphismes, then the objet representing X$\times$ is just Isom(X$\times$ Henceforth, C will be the cat\'egorie des sch\'emas with the \'etale topology (SGAD, IV, (6.3)).

\begin{quote}D\'efinition (4.5). --- A champ 50 is quasi-s\'epar\'e /f the diagonal morphisme from to\end{quote}
5O to 50 $\times$ 5O is representable, quasi-compact and s\'epar\'e.

\begin{quote}D\'efinition (4.6). --- A champ 50 is an algebraic champ (1) /f\end{quote}
(i) 5O$\to$ SOx so is representable; (ii) there exists a i-morphisme x : X \textasciitilde{} SO such that for all y : Y \textasciitilde{} SO, the projection morphisme X$\times$ is surjectif and \'etale (i.e., x is \'etale and surjectif). (1) This definition is the " right " one only for quasi-s\'epar\'e champs. It will however be sufficient for our purposes.

The 2-cat\'egorie of champs alg\'ebriques contains the representable champs and is stable under products and fibre products. If 5: is a quasi-s\'epar\'e algebraic champ, the diagonal map is unramified and quasi-affine.

\begin{quote}D\'efinition (4-7). --- An algebraic champ 5: is s\'epar\'e /f ,9\textasciitilde{} 5: is propre (or,\end{quote}
equivalently,fini). A I-morphisme f: 5:1$\to$5:2 is s\'epar\'e (resp. quasi-s\'epar\'e) if for any morphbm x : X$\to$\textasciitilde{}9\textasciitilde{} from a s\'epar\'e sch\'ema X to 5:2, the fibre product 5:1$\times$ X is s\'epar\'e (resp. quasi-s\'epar\'e). Example (4.8). --- Soit X soit un sch\'ema over S. Soit G soit un group sch\'ema over S, \'etale, s\'epar\'e and of fini type over S, which operates on X. We will denote by IX/G] the S-champ whose cat\'egorie of sections over an S-sch\'ema T is the cat\'egorie of principal homogeneous spaces (p.h.s.) E over T, with structural group G (i.e., a p.h.s, under GT), provided with a G-morphisme $\varphi$: E \textasciitilde{}X. The principal homogeneous space G $\times$ X over X (G acting only on the first factor) plus the G-morphisme G$\times$ X$\to$X (given by the action of G on X) is a section of [X/G] over X. The corresponding morphisme q : X$\to$ [X/G] is \'etale and surjectif, so that [X/G] is an algebraic champ. In addition, X is a principal homogeneous space over [X/G]; the champ [X/G] is representable if and only if X is a principal homogeneous space over a sch\'ema Y, in which case [X/G] my. If X = S, then [X/G] = [S/G] might be called the " classifying champ " of G over S. Example (4.9). --- Supposons que a champ \$f has the proprety that in each cat\'egorie \$fx the only morphismes are the identit\'e morphismes. Then 2fx is just a set o\textasciitilde{}(X) = Ob \textasciitilde{}9\textasciitilde{} and this set, under pull-back, is a contravariant foncteur o\textasciitilde{} in X. Conditions (ii) and (iii) of (4. I) assert that the foncteur o\textasciitilde{} is a sheaf on C. Artin and Knutson [K] have defined an algebraic space to soit un sheaf o \textasciitilde{} such that: (i) for any morphismes X>o\textasciitilde{}, y-$\to$o\textasciitilde{} of representable foncteurs to o\textasciitilde{}, the fibre product X$\times$ is representable; (ii) there exists a morphisme X$\to$o\textasciitilde{}, represented by surjectif, \'etale morphismes of sch\'emas. This is exactly what we have called an algebraic champ in this case.

\begin{quote}D\'efinition (4. xe). --- Soit 5\# soit unn algebraic champ. The \'etale site 5\textasciitilde{} of 5: is\end{quote}
the cat\'egorie with objets the \'etale morphismes x : X\textasciitilde{}5" and where a morphisme from (X, x) to (Y,y) is a morphisme of sch\'emas f: X$\to$Y plus a 2-morphisme between the I-morphisme x : X$\to$5 \textasciitilde{} and y.f : X$\to$S:. A collection of morphismes f\textasciitilde{} : (X\textasciitilde{}, x\textasciitilde{}) \textasciitilde{} (X, x) is a coveringfamily if the underlying family of morphismes of sch\'emas is surjectif. The site 5:et is in a natural way ringed. When we speak of sheaves on 5r we mean sheaves on Sect.

IO0 P. DELIGNE AND D. MUMFORD We now explain how many concepts from la th\'eorie des sch\'emas may soit unpplied to champs alg\'ebriques. Soit p soit un proprety of morphismes of sch\'emas, stable by \'etale change of base, and of a local nature (for the \'etale topology) on the target. For instance: being an open immersion with dense image, being dominant, birational... A representable morphisme of champs alg\'ebriques f: TI\textasciitilde{}T,, is said to have proprety P if for one (and hence for every) surjectif \'etale morphisme x:X\textasciitilde{}T2, the morphisme of sch\'emas deduced by changement de base f' : Xxr2T---\textasciitilde{}X has that proprety. Soit P soit un proprety of morphismes of sch\'emas, which, at sou$\pi$e and target, is of a local nature for the \'etale topology. This means that, for any family of commutative squares Xl gi > X 4 l, Yi hi> y where the gi (resp. hi) are \'etale and cover X (resp. Y): P(f) .r Vie(f/). For instance:f plat, lisse, \'etale, unramified, normal, localement of fini type, localement of fini presentation. If f: TI-\textasciitilde{}T\textasciitilde{} is a morphisme of champs alg\'ebriques, we say f has proprety /' if for one, then necessarily for every, commutative diagram X x > T1 4 l' y u> T 2 where X and Y are sch\'emas and x,y are \'etale and surjectif, f' has proprety/'. Similarly, if P is a proprety of sch\'emas, of a local nature for the \'etale topology, an algebraic champ T will be said to have proprety t' if for one (and hence for every) surjectif \'etale morphisme x : X$\to$T, X has proprety P. This applies to, for instance, the propreties of being regular, normal, localement noetherian, of characteristic p, r\'eduites, Cohen-Macaulay... An algebraic champ T will be called quasi-compact if there exists a surjectif \'etale morphisme x :X$\to$T with X quasi-compact. A morphisme f: TroT., of champs alg\'ebriques will be called quasi-compact if for any quasi-compact sch\'ema X over T2, the fiber product T\textasciitilde{}$\times$ is quasi-compact. It is enough to test the condition for a surjectif family fr : X$\simeq$+T2. We define a morphisme f: TI\textasciitilde{}T 2 to be of fini type, if it is quasicompact, and localement of fini type; of fini presentation, if it is quasi-compact, quasi-

\begin{quote}D\'efinition (4-II). - - A morphisme f: TI$\to$T2 is propre /f it is s\'epar\'e, of fini\end{quote}
type and if, localement over T2, there exists commutative diagrams T3 \textasciitilde{}\textasciitilde{}, T I with g surjectif and h representable and propre. The following form of Chow's Lemme will be sufficient for our purposes.

\begin{quote}Th\'eor\`eme (4-I2). --- Soit S soit un noetherian sch\'ema andf soit un morphismefrom an \'etale site T\end{quote}
to S. We assumef to be s\'epar\'e and offini type. Then, there exists a commutative diagram T ,g T' \textasciitilde{}, T" ! t' in which T' and T" are sch\'emas and such that (i) g is propre, surjectif and generically fini; (ii) j is an open immersion; (iii) f" is projective. Using (4. I2), it is easy to extend the cohomological theory of faisceaux coh\'erents to the present situation. In fact if f : TI$\to$T 2 is a propre morphisme of noetherian champs alg\'ebriques and if o\textasciitilde{} is a coherent sheaf on T\textasciitilde{}, then the R\textasciitilde{}f.($\simeq$) are faisceaux coh\'erents on T\textasciitilde{}. H\textasciicircum{}0wever, the R\textasciitilde{}f,($\simeq$) don't need to be zero for i large enough. Soit S soit un sch\'ema and G a fini group. We denote by p the projection p : [S/G]$\to$S. Quasi-coherent (resp. coherent) sheaves of modules on [S/G] may (and will) be identified with quasicoherent (resp. coherent) sheaves of modules on S, on which G acts. One has R\textasciitilde{}p,(y) \textasciitilde{}H\textasciicircum{}0(G, o\textasciitilde{}). In general, the (quasi-coherent sheaf) cohomology of champs alg\'ebriques appears as a mixture of fini group cohomology and of sch\'ema cohomology. The disjoint sum T of a family (Ti)\textasciitilde{}eI of champs is the champ a section of which over a sch\'ema X consists of (i) a decomposition X=[Ix i of X; (ii) a section of Tr over Xi for each i.

I02 P. DELIGNE AND D. MUMFORD The void champ O is the one represented by the void sch\'ema. A champ is connexes if it is non-void and is not the disjoint sum of two non-void champs.

\begin{quote}Proposition (4. x3 ). --- A localement noetherian algebraic champ is in one and only one way the\end{quote}
disjoint sum of a family of connexes champs alg\'ebriques (called its connexes composantes). We denote by n0(T) the set of connexes composantes of localement noetherian algebraic champ T. If x : X-$\to$T is surjectif and \'etale, n0(T) is the cokernel of the two maps \textasciitilde{}o(X x\textasciitilde{}X) \textasciitilde{} ,\textasciitilde{}o(X) $\to$ \textasciitilde{}o(T).

\begin{quote}Proposition (4. x4)- --- Soit T soit unn algebraic champ of fini type over a corps k. Then,\end{quote}
T is connexes if and only if there exists a connexes sch\'ema X, offini type over k, and a surjectif morphisme from X to T. An ouvert U of an algebraic champ T is a full subcat\'egorie U c T which is an algebraic champ, which contains together with any teOb(T) all isomorphe t' and such that the inclusion j:U-$\to$T is representable by open immersions. The ouverts of T corresponds bijectively to the ouverts of its \'etale site. For each ouvert U of T, there exists one and only one full subcat\'egorie T---U of T, which is an algebraic champ, which contains together with any teOb(T) all isomorphe t' and such that (i) T---U is r\'eduites; (ii) the inclusion map i:T---U-$\to$T is representable by closed immersions; (iii) for any \'etale surjectif morphisme x : X$\to$T, the image inverse of T---U on X is the complement of the image inverse of U. An algebraic champ F in T satisfying (i) and (ii) is a closed subset of T, and the foncteur U \textasciitilde{}T---U is an isomorphismee of the set of open and the set of closed subsets of T. If F satisfies only (ii), F$\pi$,\textasciitilde{}satisfies (i) and (ii) so that F defines a closed subset of T. An algebraic champ T is irr\'eductible if it is not the union of two closed subsets, non void and distinct from T.

\begin{quote}Proposition (4. x5). - - A noetherian algebraic champ T is in one and only one way the union\end{quote}
of irr\'eductible closed subsets, none of which contains any other. They are called the irr\'eductible composantes of T. If U is an open dense subset of T, the irr\'eductible composantes of U are the non-void intersections of U with the irr\'eductible composantes of T. Each irr\'eductible component Of T is contained in a connexes component of T. Conversely:

\begin{quote}Proposition (4. x6). --- The connexes composantes of a normal noetherian algebraic champ\end{quote}
are irr\'eductible.

\begin{quote}Th\'eor\`eme (4. x7). I Soit f soit un morphisme of fini type from an algebraic champ T to a\end{quote}
noetherian sch\'ema S. For seS, let n(s) be the number of connexes composantes of the fibre g\'eom\'etrique of T at s. Then (i) n(s) is a constructible function of s;

\begin{quote}Th\'eor\`eme (4. z8) (Valuative criterion for separation.) --- The morphismef is s\'epar\'e\end{quote}
if and only if, for any complete anneau de valuation discr\`ete with alg\'ebriquement clos residuecorps and any commutative diagram T Spec(V) \textasciitilde{}. S any isomorphismee between the restrictions of gl and g2 to the genericpoint o.f Spec(V) can be extended to an isomorphismee between gl and g2. This criterion is nothing other but the valuative criterion of propreness (EGA, II, 7.3 .8) applied to the (representable) diagonal morphisme.

\begin{quote}Th\'eor\`eme (4. I9) (Valuative criterion for propreness.) --- If f is s\'epar\'e, thenf is\end{quote}
propre if and only if, for any anneau de valuation discr\`ete V with corps offractions K and any commutative diagram Spec(K) , Spec(V) \textasciitilde{}. S there exists a fini extension K' of K such that g extends to Spec(V'), where V' is the cl\^oture int\'egrale of V in K' // Spec(K') ---/------\textasciitilde{} Sp.ee(V') \textbackslash\{\} Y -l \textbackslash\{\} Spec(K) > Spec(V) > S To prove a given f is propre, it suffices to verify the above criterion under the additional hypothesis that V is complete and has an alg\'ebriquement clos residue corps. Further, given a dense ouvert U of T, it is enough to test only g's which factor through U.

\begin{quote}Proposition (4.20). --- Soit 5P soit unn algebraic champ. Le foncteur which, to any algebraic\end{quote}
champ over St, f: \$'---$\to$5\textasciitilde{} associates the Ose sheaf of algebras f,O,- induces an equivalence of cat\'egories between: (i) the cat\'egorie of champs alg\'ebriques representable and affine over St'; (ii) the dual of the cat\'egorie of quasi-coherent Os$\simeq$algebras.

]04 P. DELIGNE AND D. MUMFORD Soit \textasciitilde{}r soit un quasi-coherent sheaf of d$\simeq$algebras on an algebraic champ 5\textasciitilde{} For each \'etale morphisme x : U$\to$5r with U affine, let d'(U) be the cl\^oture int\'egrale of F(U, Os\textasciitilde{}) in d(U). By (EGA, II, 6.3.4) , the d'(U) for variable U are the sections over U of a quasi-coherent sheaf d' on \textasciitilde{}9 \textasciitilde{} which will be called the cl\^oture int\'egrale of (0\textasciitilde{} in d. Soit f: \textasciitilde{}"\textasciitilde{}5r be representable and affine. The algebraic champ which is associated by (4.2o) to the cl\^oture int\'egrale of \textasciitilde{}9\textasciitilde{} in f.0$\simeq$ will be called the normalisation of 5a with respect to \$'. Its formation is compatible with any \'etale change of basis.

\begin{quote}Th\'eor\`eme (4.2I). --- Soit 5a soit un quasi-s\'epar\'e champ over a noetherian sch\'ema S.\end{quote}
Assume that (i) the diagonal map 5t'$\to$\textasciitilde{}9\textasciitilde{}215 5P is representable and unramified; (ii) there exists a sch\'ema X of fini type over S and a lisse and surjectif S-morphisme from X to 5f. Then, ,9" is an algebraic champ of fini type over S. M. Artin has developed powerful methods to relate pro-representability of a champ to the existence of \'etale surjectif maps x : X-$\to$5\textasciitilde{} § 5. Second proof of the irreducibility theorem. Soit \textasciitilde{}'g (g>2) be the champ whose cat\'egorie of sections over a sch\'ema S is the cat\'egorie of stable curves of genus g over S, the morphismes being the isomorphismees of sch\'emas over S. By (I. ii), the diagonal morphisme A :\textasciitilde{}r215 is representable, fini and unramified. We saw in w I that the champ classifying the tricanonically embedded stable curves of genus g is represented by a sch\'ema H\_g, lisse and of fini type over Spec(Z). The " forgetful " morphisme ng- g is representable, lisse and surjectif. Indeed, if p : C$\to$S is a stable curve over S, defining a morphisme \textasciitilde{}, : S$\to$\textasciitilde{}gg, then the fibre product Fig $\times$ S is the sch\'ema, lisse over S, of isomorphismees between the standard projective space of dimension 5g---6 over S and P(p,(o\textasciitilde{}c\textasciitilde{}) ). We deduce from this and (4.2 I) that:

\begin{quote}Proposition (5. x). ---..lt'g is a s\'epar\'e algebraic champ of fini type over Spec(Z).\end{quote}
Soit us denote by \textasciitilde{} the ouvert of \textasciitilde{}r which " consists of " lisse curves, and by \textasciitilde{}g the " universal curve " over dr'g, the algebraic champ classifying pointed stable curves.

\begin{quote}Th\'eor\`eme (5.2). --- The champs alg\'ebriques .ZZg and \textasciitilde{}g are propre and lisse over Spec(Z)\end{quote}
and the complement of vs \textasciitilde{} in \textasciitilde{}g is a divisor with normal crossings relative to Spec(Z).

\begin{quote}D\'efinition (5-4). --- AJacobi structure of level k on X is an isomorphismee (respecting their\end{quote}
homogeneous symplectic structures) between Rlp,(Z/kZ) and (Z/kZ)2g. (5.5) Given a section s of p, and a set of prime numbers P including all residue characteristics of S, the specialisation theorem for the fundamental group enables one to construct a pro-objet tel(X/S, s)/P/ of the cat\'egorie (l.c. gr (P)) of localement constant sheaves of fini groups of order prime to P, with the propreties (i) H\textasciicircum{}0ms(\textasciitilde{}l(X/S , s)(P), G)=R\textasciitilde{}p,(X mod s, p*G)=Rlp,(Ker(p*G\textasciitilde{}s,G)) foncteurially in Ge(1.c. gr a')) ; (ii) the formation of hi(X/S, s)(P) is compatible with any change of base.

io6 P. DELIGNE AND D. MUMFORD IfG and H are two sheaves of groups on Set, we define the sheaf of exterior morphismes from H to G, H\textasciicircum{}0me\textasciitilde{}t(H, G), as the quotient of H\textasciicircum{}0m(H, G) by the action of H induced by its action on itself by inner automorphisme. The sheaf H\textasciicircum{}0m\textasciitilde{} s)(P),G),-\textasciitilde{} R\}p,(p*G) (for G\textasciitilde{}(1.e. gr(P))) is " independent " of the choice of s. We shall denote it as Uomyt(rzl(X/S) (P),G). As p : X\textasciitilde{}S admits sections localement for the \'etale topology, this sheaf makes sense without assuming p to have a global section. It is independent of the choice of P, so long as P is prime to the order of G and includes all residue characteristics of S.

\begin{quote}D\'efinition (5.6). --- Soit G soit un fini group of order n prime to P. A Teichmfiller\end{quote}
structure of level G on X is

\displaycrop{p107_level_table_display.png}
\displaycrop{p107_level_commutative_display.png}
\displaycrop{p108_teichmuller_commutative_display.png}
\displaycrop[0.7\linewidth]{p108_exact_sequence_display.png}

 a surjectif exterior homomorphismee from nl(X/S) to G. The fini generation of z:I(X/S ) (SGA, 6o/6I, exp. IO) implies:

\begin{quote}Lemme (5.7). --- The sheaf on (Sch/S) of the Teichm\"uller structures of level G on X is\end{quote}
represented by an \'etale covering of S. We denote by \textasciitilde{},0 the champ classifying the stable lisse curves of genus g and G g characteristic prime to n, with a Teichmfiller structure of level G. For any algebraic champ dr', we denote by \textasciitilde{}s its ouvert dr'$\times$ Spec(Z[i/n]).

\begin{quote}Lemme (5.7) may now be rephrased:\end{quote}
\begin{quote}Proposition (5.8).- The " forgetful " morphisme\end{quote}
is representable, fini and \'etale. The champ a.A'\textasciitilde{}thus in an algebraic champ. Soit e\textasciitilde{}go be the normalisation of.g 0[I In] with respect to a.g \textasciitilde{} The champ o.g0, being representable and fini over .g0[I/n], is propre over Spec(Z[i/n]).

\begin{quote}Th\'eor\`eme (5.9). --- The geometricfibres of the projection of a.go onto Spee(Z [I/n]) are\end{quote}
normal, and, fibre by fibre, \textasciitilde{}.g\textasciitilde{} is dense in \textasciitilde{}.go" D\'emonstration. --- We will use Abhyankar-Artin's lemma, in its " absolute " form:

\begin{quote}Lemme (5. xe). --- Soit D soit un divisor with normal crossings on an excellent regular\end{quote}
sch\'ema X, Y an \'etale covering of X---D and Y the normalisation OfX with respect to Y. Assume that the generic points of the irr\'eductible composantes of D are all of characteristic o. Then every geometric point of X has an \'etale neighbourhood x : X'$\to$X such that, on X': (i) D becomes a union of regular divisors (D,),\textasciitilde{}I, D, of equation t$\simeq$\textasciitilde{}o. (ii) There exists an integer k prime to residue characteristics of X' such that Y becomes isomorphe to a disjoint union of quotients (by subgroups of gi) of the covering of X' obtained by extracting the Uh-roots of the ti's. If x: X$\to$.gg[ I] is an \'etale morphisme, then X\textasciitilde{} is the complement L---..I in X of a divisor with normal crossing relative to Spec(Z), X $\simeq$---=X $\times$ (aMl\textasciitilde{}g) is an \'etale

\begin{quote}Corollaire (5. xI ). --- The fibres g\'eom\'etriques of the projection\end{quote}
G g $\to$ Spec Z all have the same number of connexes composantes. D\'emonstration. --- By (4.17), all fibres g\'eom\'etriques odt'gX Spec(Fz) of a.\textasciitilde{}t'g over Spec(Z[I/n]) have the same number of connexes composantes. These connexes composantes are irr\'eductible (4- 16). Furthermore a\textasciitilde{}'$\Omega$\_XSpec(Fz) is dense in adt'g$\times$ Spec(F\textasciitilde{}) and thus has the same number of connexes (or irr\'eductible) composantes as Q.A\textasciitilde{}'g X Spec(Fz), and this number is independent of I. (5- x2) Soit us denote by H the fundamental group of an oriented closed differentiable surface SOof genus g. The group II may be defined by generators and relations as follows: (i) generators: xi for i < i < 2g; (ii) relation: (x\textasciitilde{},xa+l). . . (x\textasciitilde{}, xa+\textasciitilde{}). . . (xg, x2a)=e , where (a, b)=aba-lb -1. Soit (e\textasciitilde{})be the standard basis of Z2g. The morphisme $\varphi$from H to Zegwith $\varphi$(Xi) = el identifies H/(H, n)=Hl(n, Z) with Z\textasciitilde{}a. The surface SOis a K(II, i) and thus H\textasciicircum{}0(H, Z)= I-P(S0, Z)=Z, and the cup product defines a symplectic structure on II/(II, H). This structure is identified by $\varphi$ with the standard symplectic structure of Z2g. We denote by Aut\textasciitilde{} the subgroup of Aut(II) which acts trivially on H2(II, Z). Dehn has proved that each exterior automorphisme of II is induced by a diffeomorphisme of SOonto itself and that the map induced by $\varphi$: Aut\textasciitilde{} -$\to$ Sp\textasciitilde{}(Z), is surjectif (see[Ma]).

\begin{quote}Th\'eor\`eme (5. I3). --- The number of connexes composantes of any fibre g\'eom\'etrique of the\end{quote}
projection of Q\textasciitilde{}o onto Spec(Z[I/n]) is equal to the number of orbits of Aut\textasciitilde{} in the set of exterior epimorphismes from H to G. D\'emonstration. By (5. I3), it suffices to prove that 0 --- Gdt'gxSpec(C) has the said number of connexes composantes. As results from (4-14), rr0(a-/t'\textasciitilde{} X Spec(C)) = r:0(GM\_g) where aM\_g is the coarse moduli sch\'ema classifying stable lisse curves of genus g over C with a Teichm\"uller structure of level G.

\textasciitilde{}o8 P. DELIGNE AND D. MUMFORD Recall that a Teichm\"ulle$\pi$urveof genus g is a stable lisse curve C of genus g over (I together with an exterior isomorphismee $\varphi$of the transcendental fundamental group rq(C) with H, which induces a symplectic isomorphismee (1) between HdC , Z) \textasciitilde{} \textasciitilde{}dC)/(\textasciitilde{}dC), \textasciitilde{}1(C)) and II/(II, H). By Teichmfiller's theory [W], the analytic space T0 classifying Teichm\"uller curves of a given genus g>2 is homeomorphic to a ball, and hence connexes. If \textasciitilde{}$\varphi$is a surjectif homomorphismee from II to G, the map (C, 9) \textasciitilde{} (C, +9) defines a morphisme t+ : To-$\to$GM0. Two such maps t+ and t+, have the same image if and only if + = +'\textasciitilde{} for \textasciitilde{}Aut\textasciitilde{} and H t\textasciitilde{},(T0) GMa= r mod Aut\textasciitilde{} which implies (5- i3). (5. x4) Soit us denote by n\textasciitilde{}'\textasciitilde{} the algebraic champ classifying stable lisse curves with a Jacobi structure of level n. This algebraic champ " is" a true sch\'ema for n> 3 (by [S]). If 9 is a Jacobi structure of level n on a stable lisse curve p : X \textasciitilde{}S, we define the " multiplicator " \textasciitilde{}(9) of 9 by the commutative diagram A (z/nz)20 z/nz A2Rlp,(g/nZ) A> \textasciitilde{}\_\textasciitilde{} K2p,(g/ng) The sch\'ema of isomorphismees between Z/nZ and $\simeq$1 is Spec(Z[e\textasciitilde{}'\textasciitilde{}j", i/n]) thus 9\textasciitilde{}[\textasciitilde{}(9) defines a morphisme \textasciitilde{} from S to Spec(Z[e2\textasciitilde{}/", I/n]). This being true for any X and S, [z is induced by : ,dr'\textasciitilde{}\textasciitilde{} Spec(Z[e\textasciitilde{}/n, I/n]).

\begin{quote}Th\'eor\`eme (5. x5). --- The geometricfibres of the morphisme \textasciitilde{} are connexes.\end{quote}
\begin{quote}D\'emonstration. \textasciitilde{} By definition, n\textasciitilde{}'\textasciitilde{} is open and closed in G Jr'0 \textasciitilde{} for G=(Z/nZ) \textasciitilde{}. The group GL2g(Z/nZ) acts on G, and thus on odt'\textasciitilde{} One has: (i) the ouvert ,all\textasciitilde{} of \textasciitilde{}/t'\textasciitilde{} is stable under the subgroup H = CSp2g(Z/nZ ) of symplectic similitudes;\end{quote}
(ii) -J = II o( de' tl 0 G/H \textasciitilde{}n gJ 9 (1) An arbitrary isomorphismee 9 induces an isomorphismee betwen Hx(C, Z) and II/(II, II) which is always symplectic up to sign.

\begin{quote}Lemme (5- x6). --- The homomorphismee\end{quote}
Aut\textasciitilde{} $\to$ Sp\textasciitilde{}(Z/nZ) is surjectif. D\'emonstration. --- This results from Dehn's theory (5. I2) and from the fact that the groups Sp,, being split semi-simple simply connexes groups, are generated by their unipotent elements, so that the reduction map sp /z) Sp /Z/ Z) is surjectif. BIBLIOGRAPHY [Ab] S. ABHYANKAR, Resolution of singularities of arithmetical surfaces, Proc.Conf. on Arith. Alg. Geom. at Purdue. 1963. [A1] M. ARTIN, The implicit function theorem in algebraic geometry, Proc. Colloq. Alg. Geom., Bombay, i968. [As] M. ARTIN, Some numerical criteria for contractibility, Am. 07. Math., 84 (I96\textasciitilde{}), p. 485. [B] J. BENABOO,Th6se, Paris, I966. [E] F. ENRIO.UESand CH\textasciicircum{}1SlNI, Teoria geometrica delle equazioni e deUefunzioni algebriche, Bologna, 1918. IF] W. FULTON,Hurwitz sch\'emas and irreducibility of moduli of algebraic curves, Ann. of Math. (forthcoming). [G] J. GIRAUD,Gohomologienon abdlienne, University of Columbia. [Gr] A. GROTI-IENDmCK,Techniques de descente et th6or6mes d'existence en g6om6trie algdbrique, IV, Sem. Bourbaki, 9.21, I96O-I96I. [H] R. HARTSHORNE,Residues and Duality, Springer Lecture Notes, 20, i966. [K] D. KNUTSON,Algebraic spaces, Thesis, M.I.T., 1968. [Li] J. LI1,M\_\textasciitilde{}N,Rational singularities, Publ. Math. LH.E.S., n\textasciitilde{} 36 (I969). [L] M. LICI-IT\textasciitilde{}NBAV\textasciitilde{},Curves over anneau de valuation discr\`etes, Am. 07. Math., 90 (1968), p. 38o. [Ma] W. MANGm\%Die Klassen yon topologisehen Abbildungen einer geschlossenen Flfiche auf sich, Math. geit., 44 (1939) , p. 54I. [M1] D. MUmfORD, GeometricInvariant Theory, Springer-Verlag, I965. [1VIii D. 1ViI.$\pi$ORD, Notes on seminar on moduli problems, Supplementary mimeographed notes printed at A.M.S. Woods H\textasciicircum{}0le Summer Institute, I964. [Ms] D. MVM\_gORD,Picard groups ofmoduli problems, Proe. Conf. on Arith. Alg. Geom.at Purdue, 1963. [N] A. N\textasciitilde{}RON, Mod\&les minimaux des varidt6s ab61iennes sur les corps locaux et globaux, Publ. Math. LH.E.S., n \textasciitilde{} 31. JR] M. RAyNAUD, Spdcialisation du foncteur de Picard, Comptesrendus Acad. Soi., 264, p. 941 and p. 1ooi. [\textasciitilde{}] I. \textasciitilde{}AFAm\textasciitilde{}VlCH,Lectures on minimal models, Tata Institute Lectures notes, Bombay, I966. [So] M. SCHLESSINOER, Thesis, Harvard. IS] J.-P. SERRE, Rigiditd dufoneteur de 07aeobidYehelon n\textasciitilde{}> 3, app. A l'expos6 17 du s6minaire Caftan 6o/6I. [Se] F. SEWRI and L6FFLER, Vorlesungen iiber algebraische Geometrie, Teubner, I924. [W] A. WELL, Modules des surfaces de Riemann, Sdm. Bourbaki, 168, 1957-58. Manuscrit refu le 21 janvier 1969.

\end{document}
